\pdfoutput=1
\documentclass[11pt,a4paper]{article}

\usepackage[T1]{fontenc}
\usepackage[utf8]{inputenc}
\usepackage{amsmath,amssymb,mathtools}
\usepackage{amsthm}
\usepackage{txfonts}
\usepackage{booktabs}
\usepackage{array}
\usepackage{longtable}
\usepackage[margin=2.4cm]{geometry}
\usepackage{tikz}
\usetikzlibrary{cd}
\usetikzlibrary{decorations.pathmorphing}
\usepackage{xcolor}
\usepackage[colorlinks=true,linkcolor=blue!50!black,citecolor=blue!50!black,urlcolor=blue!50!black]{hyperref}

\theoremstyle{plain}
\newtheorem{theorem}{Theorem}[section]
\newtheorem{proposition}[theorem]{Proposition}
\newtheorem{lemma}[theorem]{Lemma}
\newtheorem{corollary}[theorem]{Corollary}
\newtheorem{conjecture}[theorem]{Conjecture}
\theoremstyle{definition}
\newtheorem{definition}[theorem]{Definition}
\newtheorem{example}[theorem]{Example}
\newtheorem{task}{Problem}
\theoremstyle{remark}
\newtheorem{remark}[theorem]{Remark}
\newtheorem{caveat}[theorem]{Remark}
\newtheorem*{remarkun}{Remark}

\newcommand{\cat}[1]{\mathsf{#1}}
\newcommand{\catU}{\cat{U}}
\newcommand{\catV}{\cat{V}}
\newcommand{\catO}{\cat{\Omega}}
\newcommand{\catS}{\cat{S}}
\newcommand{\catD}{\cat{D}}
\newcommand{\catK}{\cat{K}}
\newcommand{\Meas}{\cat{Meas}}
\newcommand{\Set}{\cat{Set}}
\newcommand{\Stoch}{\cat{Stoch}}
\newcommand{\FinStoch}{\cat{FinStoch}}
\newcommand{\FinSet}{\cat{FinSet}}
\newcommand{\Alg}{\cat{Alg}}
\newcommand{\G}{\mathcal{G}}
\newcommand{\Prb}{\mathcal{P}}
\newcommand{\op}{\mathrm{op}}
\newcommand{\id}{\mathrm{id}}
\newcommand{\pr}{\mathrm{pr}}
\newcommand{\del}{\delta}
\newcommand{\R}{\mathbb{R}}
\newcommand{\N}{\mathbb{N}}
\newcommand{\Vect}{\cat{Vect}}
\newcommand{\Norm}{\mathcal{N}}
\newcommand{\Bern}{\mathrm{Bern}}
\newcommand{\expit}{\mathrm{expit}}
\newcommand{\res}{\mathrm{res}}
\DeclareMathOperator{\ob}{ob}
\DeclareMathOperator{\Nat}{Nat}
\DeclareMathOperator{\Hom}{Hom}
\DeclareMathOperator{\HH}{HH}
\newcommand{\Cat}{\cat{Cat}}
\newcommand{\StatMod}{\cat{StatMod}}
\newcommand{\Statk}{\StatMod_{\Meas}}
\newcommand{\Jc}{\Delta}            % comparison functor StatMod -> Arrdet(Stoch)
\newcommand{\Jcm}{\Delta_{\Meas}}
\newcommand{\Arrdet}[1]{#1^{\rightarrow}_{\det}}
\newcommand{\Arrcop}[1]{#1^{\rightarrow}_{\det,\,\mathrm{cop}}}
\newcommand{\Arrim}[1]{#1^{\rightarrow}_{\del}}
\newcommand{\fnote}[1]{\textcolor{blue}{FV: #1}}

\title{Statistical models as natural transformations:\\
a bridge from the McCullagh--Br\o ns hexagon to Markov categories}
\author{Francesco Vaccarino\\
\small Dipartimento di Scienze Matematiche ``G.\ L.\ Lagrange'', Politecnico di Torino}
\date{August 2026}

\begin{document}
\maketitle

\begin{abstract}
\noindent
McCullagh's 2002 categorical definition of a statistical model, sharpened by
Br\o ns into a hexagonal diagram of categories and functors, identified
naturality as the mathematical content of the informal requirement that the
meaning of a parameter be ``retained'' across experimental designs. 
The programme
supplied a language but not a calculus: no ambient categorical probability
theory existed in which the naturality condition could be computed with,
composed, or mechanically verified. We show that the Markov-categorical
framework of Fritz, Cho--Jacobs and Perrone supplies exactly this calculus. Our
main result recasts a McCullagh--Br\o ns model as a natural family of morphisms
in the Kleisli category of the Giry monad, deterministic on the parameter side;
under this translation McCullagh's criterion for a meaningful subparameter
becomes a naturality condition, hence a type-level property under the functorial
encoding, and Tjur's weaker elementary criterion, invariance under design
reduction, is recovered as naturality restricted to the insertion morphisms.
We prove the statement in full for finite (discrete) models, state the
measurable case with its precise additional hypotheses, and formulate the
coherence problem for finite designs in the cohomology of a diagram of algebras
in the sense of Gerstenhaber--Schack and Baues--Wirsching, connecting it to
recent work of Caputi and the author. We show that the obstruction so obtained
is necessarily \emph{class-relative}, the unrestricted class vanishes for
every quantity, by construction, and discharge the resulting statement by an
exact computation for the one-way layout, which certifies the non-meaningfulness of
the marginal dispersion relative to the class of information-bearing
recalibrations for scale-confounded designs, recovers the within-group
dispersion as a correction by pure linear algebra, and
exhibits $H^{1}$ of the design scheme as a computable invariant of independent
interest.
\end{abstract}

\medskip
\noindent\emph{MSC 2020:} 62A01 (primary); 18M05, 60A05, 18G90 (secondary).\\
\emph{Keywords:} statistical model; Markov category; Giry monad; natural
transformation; meaningfulness; design category; Baues--Wirsching cohomology.


\tableofcontents

\section{Introduction}

A recurring observation is that a small canon of unifying principles holds statistics together:
likelihood, sufficiency, conditioning, ancillarity, invariance,
exchangeability. Each principle is made precise inside its own apparatus
(measure-theoretic probability, decision theory, group actions), and their
mutual coherence is maintained by the trained intuition of the field rather
than by a single calculus in which they are theorems. The best-known sustained
attempt at such a calculus is Birnbaum's~\cite{birnbaum_foundations_1962}, which
formulates sufficiency and conditionality as axioms on evidential meaning and
argues that together they force the likelihood principle. That its status
remains disputed sixty years on, over what the conditionality principle
\emph{says}, rather than over any computation~\cite{mayo_birnbaum_2014,
evans_birnbaum_2013}, is itself evidence that the principles have lacked an
ambient formalism in which such a question is settled. McCullagh's ``What is a
statistical model?''~\cite{mccullagh_what_2002} diagnosed one root of this state of
affairs: the standard definition of a statistical model, a parametrised family
of probability measures on a fixed sample space, is silent about how the family
must transform when the experimental design changes, and therefore cannot even
\emph{state}, let alone enforce, the requirement that a parameter retain
its meaning across designs. His remedy was categorical: designs form a
category, sample spaces and parameter spaces become functors on it, and a model
is a natural transformation. Br\o ns, in his discussion of the
paper~\cite[pp.~1280--1282]{brons_discussion_2002}, organised these data into a
hexagonal diagram of categories and connecting functors, and observed that the resulting notion of
model is equivalent, by a general comma-category argument, to a diagram in a
fixed ``category of statistical models''. In the same passage he records that
this category is too small for the purpose, to include all statistical
transformations, he writes, ``Markov kernels must be allowed'', and then
sets the point aside~\cite[p.~1282]{brons_discussion_2002}. The present paper
takes it up. Tjur, in the same discussion
volume~\cite{tjur_discussion_2002}, isolated the operative consequence, and did so
\emph{without} categorical language, which he judged convenient rather than
necessary: fixing an infinite population of covariate--response pairs, he called
a rule assigning a parameter, a function of the unknown distribution,
to each finite sample of covariate values, repetitions allowed, a
\emph{parameter function}, and such a function \emph{meaningful} when it is
invariant under the formation of marginal models, when the design is
reduced by removing some of the sampled $x$'s, the parameter of the marginal
distribution must agree with the parameter for the bigger sample. The
removals are of sample points, not of points of a covariate space: before
the criterion is imposed, his parameter functions include the sample size
$n$ itself.
His non-example is the design-averaged expectation $\alpha + \beta\bar x$,
with $\bar x$ the average over the sample one happened to draw; he attributes the
analogous failure for the correlation coefficient to McCullagh. The
categorical form of this criterion, naturality, is McCullagh's
own~\cite[\S4.5]{mccullagh_what_2002}, his rejoinder reads Tjur's remarks as
explaining what is meant by a natural parameter in
regression~\cite[p.~1304]{mccullagh_what_2002}: a \emph{natural subparameter} is a natural
transformation of functors on the design category, and inference is meaningful
only for natural parameters. The two formulations are stated over different
morphism classes, Tjur's over design reduction alone, McCullagh's over the
full design category, and neither source asserts that they coincide as
mathematical requirements.

The programme is widely admired and rarely used. We contend that the reason is
structural rather than sociological: in 2002 the codomain of the construction
was bare $\Set$, so naturality could be written down, but nothing could be
\emph{done} with it, there was no ambient categorical probability theory in
which conditioning, sufficiency, or almost-sure reasoning were native
operations and no technology for mechanically verifying naturality. Both
missing ingredients have since appeared, independently and for unrelated
reasons. Markov categories~\cite{fritz_synthetic_2020,cho_disintegration_2019,perrone_markov_2024,
fritz_definetti_2021} turn Markov kernels into the morphisms of a symmetric
monoidal category with enough structure to state and prove, synthetically, the
Fisher--Neyman factorisation, versions of sufficiency and ancillarity, and
de~Finetti-type theorems.

This paper does the obvious-in-retrospect thing: it moves the codomain of the
McCullagh--Br\o ns construction from $\Set$ to the Markov category $\Stoch$
(finite case: $\FinStoch$) and records what becomes true, what becomes
provable, and what becomes computable. The contributions are deliberately
modest in scope and precise in status:


\begin{enumerate}
\item Section~\ref{sec:mcb} gives a self-contained account of the
McCullagh--Br\o ns framework, including a complete statement and proof of the
comma-category equivalence (Proposition~\ref{prop:brons}), which in the
original discussion is asserted with a sketch.
\item Section~\ref{sec:markov} recalls the Markov-categorical background
actually used.
\item Section~\ref{sec:bridge} states and proves the bridge in the finite case
(Theorem~\ref{thm:bridge-finite}), states the measurable case with its exact
additional hypotheses (Theorem~\ref{thm:bridge-meas}), and is explicit about
the fact that those hypotheses are a genuine restriction relative to
McCullagh's $\Set$-valued formulation (Remark~\ref{rem:hypotheses}). It also
records, opening \S\ref{sec:Bref}, that the bridge is, in disguise, a representability-plus-copower
computation (Remark~\ref{rem:repr}), and uses this to
reformulate what we call the \emph{arrow-category identification}: the open
question whether Br\o ns' category of statistical models is equivalent to a
category built canonically from $\Stoch$. The answer is that, on the objects
whose parameter set and sample space satisfy two explicit regularity
hypotheses (automatic in the finite case), it is equivalent to the arrows of
$\Stoch$ \emph{with copower domain} (Proposition~\ref{wprop:copower}); for the
measurable comma category the same holds without the copower restriction
(Remark~\ref{cav:Bscope}); the passage
between the two is an adjunction exhibiting Br\o ns' category as a coreflective
subcategory of its measurable analogue, so that the copower restriction is the
non-invertibility of a counit and nothing more
(Proposition~\ref{prop:coreflection} and
Corollary~\ref{cor:coreflection-diagrams}).
\item Section~\ref{sec:meaningful} defines Tjur meaningfulness precisely,
proves its characterisation as a naturality condition under the bridge
(Proposition~\ref{prop:tjur}), and reduces its verification over finite
designs to finitely many matrix equalities (Corollary~\ref{cor:decidable}).
\item Section~\ref{sec:hochschild} formulates the coherence problem for finite
designs cohomologically: a candidate design-indexed quantity is correctable,
within a prescribed admissible class, to a natural (hence meaningful) one iff an
explicit obstruction cochain in the cohomology of a diagram of algebras \`a la
Gerstenhaber--Schack~\cite{gerstenhaber_schack_1983}/Baues--Wirsching~\cite{baues_cohomology_1985}
is the coboundary of a correction in that class.
The equivalence is proved (Proposition~\ref{wprop:obstruction}),
together with the observation that forces the obstruction to be read relative to
an admissible correction class. The one-way layout is then carried out in full
(Examples~\ref{ex:oneway}--\ref{ex:oneway-h1}).
\item Section~\ref{sec:nonclaims} delimits the scope. Several statements
adjacent to the results below (paradigm unification, nonparametric limits)
are formulable in this framework but are not proved here; the boundary is
drawn explicitly rather than left blurred.
\end{enumerate}

\paragraph{Related work.}
The closest prior work is Patterson's thesis~\cite{patterson_algebra_2020}, which
develops an algebraic theory of statistical models with machine representation
in view; the present paper differs in taking the design-indexed
(McCullagh--Br\o ns) structure, rather than the model-theoretic signature, as
primitive and specifically targeting the naturality/meaningfulness axis.
Fritz~\cite{fritz_synthetic_2020} proves synthetic sufficiency theorems for a single
fixed sample space, \emph{inside} the hom-sets of one Markov
category~\cite[Def.~14.1]{fritz_synthetic_2020}; McCullagh--Br\o ns naturality is a
condition \emph{across} a diagram of such hom-sets; the two are complementary,
and their combination (fibred over the design category) is exactly the
structure studied here. Jacobs' channel-based
probability~\cite{jacobs_structured_2021}, the synthetic de Finetti
theorem~\cite{fritz_definetti_2021}, the representable-category comparison of
statistical experiments~\cite{fritz_representable_2023} and the dilational
axiomatics of~\cite{fritz_dilations_2023} supply techniques we expect to
import in later work. Categories of Markov kernels have also been put to work
well outside inference proper: Sergeant-Perthuis develops rigorous statistical
mechanics compositionally, with statistical systems presented as
representations of posets and their Gibbs phases as invariants of those
representations, in a setting that couples measurable maps with Markov
kernels~\cite{sergeantperthuis_categorical_2023,
sergeantperthuis_compositional_2024}. That the same category of kernels underlies
both that setting and the design-indexed one studied here is evidence of its
generality:
the two settings are unrelated here, and no correspondence is claimed.
Helland's work on
statistical symmetry~\cite{helland_steps_2010} pursues the invariance principle by
group-theoretic rather than categorical means; the compositional approach to
learning of~\cite{fong_backprop_2019} is adjacent to, but distinct from, the
design-indexed structure studied here.

Two further observations in Br\o ns' discussion bear on the shape of what
follows, and we record them here because they are easily mistaken for
departures from the 2002 framework when they are in fact requests made within
it. First, having recognised a model as a diagram in his category of
statistical models over the basic category, he remarks that no evident reason
confines the theory to base categories of that particular form, and proposes
admitting diagrams over arbitrary categories, which would in particular
subsume the passage to subcategories~\cite[p.~1282]{brons_discussion_2002}.
Extending the indexing beyond the comma-category shape of
Definition~\ref{def:designdata} is therefore continuation rather than
importation. Second, he observes that the product and repetition
constructions of~\cite{mccullagh_what_2002} are best treated with monoidal
categories, monoidal functors and monoidal natural transformations, noting
that $\Prb$ carries a monoidal structure given by the product of
measures~\cite[pp.~1282--1283]{brons_discussion_2002}. The codomain adopted
here is monoidal in exactly that sense; the copy--discard structure recalled
in \S\ref{sec:markov} is a strengthening of it, and is what makes conditioning
and independence expressible rather than merely the product of models.

\section{The McCullagh--Br\o ns framework}\label{sec:mcb}

We fix the notation of Br\o ns' discussion~\cite{brons_discussion_2002}, in the form used
throughout the sequel. All categories in this section are ordinary (locally
small) categories.

\begin{definition}[Design data]\label{def:designdata}
The framework is parametrised by:
\begin{enumerate}
\item a category $\catU$ of \emph{statistical units} (typical objects: finite
sets of experimental units; typical morphisms: injections, or restriction to
subsamples);
\item a category $\catO$ of \emph{covariate spaces};
\item a category $\catV$ of \emph{response scales} (typical morphisms:
measurable surjections implementing coarsening of the response, affine
rescalings, etc.);
\item the category $\catS \coloneqq \catV \times \catU^{\op}$ of \emph{sample
space indices}, together with a functor $\Gamma \colon \catS \to \Meas$ sending
$(v,u)$ to the measurable space of $v$-valued responses on the units $u$
(canonically $\Gamma(v,u) = v^{u}$ with product structure when $v$ is a measurable
space and $u$ a finite set);
\item the category of \emph{designs}
\[
\catD \;\coloneqq\; (J_{\catU} \downarrow J_{\catO}),
\]
the comma category of a pair of functors $J_{\catU}\colon \catU \to \cat{C}$,
$J_{\catO} \colon \catO \to \cat{C}$ into a common ambient category (in the
basic examples $\cat{C} = \Set$ and both $J$'s are the underlying-set
functors). An object of $\catD$ is thus a triple $(u, \omega, x)$ with
$x \colon J_\catU u \to J_\catO \omega$ a \emph{design map} assigning
covariate values to units; a morphism
$(u,\omega,x) \to (u',\omega',x')$ is a pair of morphisms
$\varphi \colon u \to u'$ in $\catU$, $\psi \colon \omega \to \omega'$ in
$\catO$ with $x' \circ J\varphi = J\psi \circ x$. We write
$\pr_{\catU} \colon \catD \to \catU$ and
$\pr_{\catO} \colon \catD \to \catO$ for the projections;
\item a concrete category $\catK$ of \emph{parameter objects}, with faithful
functor $U_{\catK} \colon \catK \to \Set$, and a contravariant \emph{parameter
functor} $\Theta \colon \catO^{\op} \to \catK$.
\end{enumerate}
\end{definition}

\begin{remark}
Contravariance of $\Theta$ encodes the statistical fact that enlarging the
covariate space enlarges (or reparametrises) the parameter, while a morphism
of covariate spaces induces a map of parameters in the opposite direction; the
regression examples of~\cite[\S6]{mccullagh_what_2002} are of this shape.
Contravariance of the $\catU$ factor in $\catS$ encodes that restricting to a
subsample induces a (projection) map of sample spaces in the opposite
direction.
\end{remark}

Let $\Prb \colon \Meas \to \Set$ denote the functor sending a measurable space
to its set of probability measures and a measurable map to the pushforward.
(This is the underlying-set composite of the Giry monad
$\G \colon \Meas \to \Meas$~\cite{giry_categorical_1982}.) Note that $\Prb$ is
$\Set$-valued: a map $\Theta_0 \to \Prb(S)$ is a family of probability
measures on $S$ indexed by $\Theta_0$, with no measurability imposed in the
index. This is the gap the bridge of \S\ref{sec:bridge} has to close,
vanishing when the parameter object is discrete, and otherwise an added
hypothesis (Theorem~\ref{thm:bridge-meas}).

\begin{definition}[McCullagh--Br\o ns model]\label{def:mcbmodel}
With the data of Definition~\ref{def:designdata}, form the two composite
functors $\catV \times \catD^{\op} \to \Set$:
\[
L \;\coloneqq\; U_{\catK} \circ \Theta \circ \pr_{\catO}^{\op}
\circ \pr_{2}
\qquad\text{and}\qquad
R \;\coloneqq\; \Prb \circ \Gamma \circ \bigl(\id_{\catV} \times \pr_{\catU}^{\op}\bigr),
\]
where $\pr_2 \colon \catV \times \catD^{\op} \to \catD^{\op}$ is the
projection. A \emph{statistical model} is a natural transformation
\[
P \colon L \Longrightarrow R.
\]
Concretely: for every response scale $v$ and every design $d = (u,\omega,x)$,
a map
\[
P_{v,d} \colon \; U_\catK\Theta(\omega) \longrightarrow \Prb\bigl(\Gamma(v,u)\bigr),
\qquad \theta \longmapsto P_{v,d}(\cdot \mid \theta),
\]
assigning to each parameter value a probability measure on the sample space,
such that for every morphism $(\alpha,(\varphi,\psi)^{\op})$ of
$\catV \times \catD^{\op}$ the evident square commutes.
\end{definition}

The six categories $\catU,\catV,\catO,\catS,\catD,\catK$ and the functors
between them, with $\Meas$ and $\Set$ as ambient codomains, assemble into
Br\o ns' hexagonal diagram; a model is a natural transformation filling the
hexagon between its two composite paths:
\[
\begin{tikzcd}[column sep=1.6em, row sep=1.9em]
& \catV \times \catD^{\op}
\arrow[dl, "\pr_2"'] \arrow[dr, "\id_\catV \times \pr_\catU^{\op}"] & \\
\catD^{\op} \arrow[d, "\pr_\catO^{\op}"'] & & \catS \arrow[d, "\Gamma"] \\
\catO^{\op} \arrow[d, "\Theta"'] \arrow[rr, phantom, "\overset{P}{\Longrightarrow}" description] & & \Meas \arrow[d, "\Prb"] \\
\catK \arrow[dr, "U_\catK"'] 
& & \Set \arrow[dl, "id_\Set"]\\
&\Set
\end{tikzcd}
\]

\begin{remark}[Br\o ns' packaging and McCullagh's own definition]
The hexagon is Br\o ns'. McCullagh's definition
in~\cite[\S4.2]{mccullagh_what_2002} holds the response scale \emph{fixed}
(\S4.7 of the same paper lets it vary): with the scale fixed, the parameter is a
contravariant functor $\catO \to \catK$, and a model is a natural transformation
of functors on $\catD$ alone. The product categories
$\catS = \catV \times \catU^{\op}$ and
$\catV \times \catD^{\op}$ are McCullagh's~\cite[\S4.7, p.~1240]{mccullagh_what_2002};
Br\o ns restates them in his own notation~\cite[p.~1280]{brons_discussion_2002}
and then packages $\Theta$ as a functor on $\catO^{\op}$ alone. Three differences are worth recording, since we use
Br\o ns' formulation throughout.

\begin{enumerate}
\item[(i)] \emph{The $\catV$ slot is additional data.} In the hexagon $\catV$
varies while $\Theta$ is defined on $\catO^{\op}$ alone, so a non-identity
response-scale morphism constrains the model at a fixed parameter value. In
McCullagh's own worked construction the parameter of a linear model is a
subrepresentation $\Theta_{\omega} \subseteq v^{\omega}$ of the standard
representation~\cite[p.~1235]{mccullagh_what_2002}, and therefore \emph{does}
depend on the response scale. The two agree exactly when $v$ is held fixed,
which is the hypothesis of McCullagh's \S\S4.2--4.6; the discrepancy is otherwise real,
and is what Remark~\ref{cav:scaleV} measures.
\item[(ii)] \emph{Morphism restrictions.} McCullagh takes the morphisms of
$\catU$ and $\catO$ to be injective~\cite[\S4.1]{mccullagh_what_2002}; Br\o ns
judges such restrictions dispensable and does not impose
them~\cite[p.~1281]{brons_discussion_2002}. Definition~\ref{def:designdata}
follows Br\o ns for $\catO$; for $\catU$, every example below takes
$\FinSet_{\mathrm{inj}}$, and the verifications of \S\ref{sec:worked} do use
the injectivity of unit maps (see the remark following~\eqref{eq:commacond}),
so that restriction is not dispensable there. The choice for $\catO$ is not
cosmetic: it fixes the class over which
naturality is quantified, and hence the gap between McCullagh's meaningfulness
criterion and Tjur's, as \S\ref{sec:meaningful} makes precise.
\item[(iii)] \emph{Nothing below depends on the difference.} Every statement
in the sequel is a statement about the naturality squares of
Definition~\ref{def:mcbmodel}, and for fixed $v$ these are exactly McCullagh's
squares of~\cite[p.~1236]{mccullagh_what_2002}. Where the $\catV$ slot does
work of its own, we say so.
\end{enumerate}
\end{remark}

\begin{example}[One-way layout]\label{ex:oneway-setup}
\leavevmode
This is the running example of the paper. We give it in the same detail as the
instantiations of \S\ref{sec:worked}, which it precedes only because it is the
simplest of the three; the verification of naturality is given at the end of
\S\ref{sec:worked}, once the naturality squares have been displayed.

\medskip
\noindent\textbf{The experiment, in words.} A single \emph{factor} is applied
to each experimental unit: every unit receives exactly one of finitely many
mutually exclusive interventions, its \emph{treatment}. There is no blocking,
no repeated measurement, and no explanatory variable other than the treatment
received, this is what \emph{one-way} means. A \emph{treatment label set}
is a finite non-empty set $\omega$ whose elements name the available
treatments (equivalently, the \emph{levels} of the factor). Only the naming is
retained: $\omega$ carries no order, no metric and no algebraic structure.
That is precisely why the covariate category below consists of bare finite
sets rather than of vector spaces as in Example~\ref{ex:linreg} below; the passage
between the two descriptions is dummy coding, and is
Remark~\ref{rem:dummy}.

\begin{enumerate}
\item \textbf{Units.} $\catU = \FinSet_{\mathrm{inj}}$: an object $u$ is an
index set for the experimental units actually enrolled; a morphism
$\varphi \colon u \hookrightarrow u'$ is the inclusion of a subsample. We
write $n = |u|$.

\item \textbf{Covariate spaces $=$ treatment label sets.} $\catO =
\catO_{\mathrm{lab}}$, the full subcategory of $\FinSet$ on the finite
non-empty sets: objects are treatment label sets, morphisms are \emph{all}
maps $\psi \colon \omega \to \omega'$ of such sets. Four kinds of morphisms
carry the statistical readings, and every $\psi$ is a surjection followed by
an injection, so the first two generate:
\begin{itemize}
\item $\psi$ \emph{surjective}: \emph{merging} of treatments. The fibres
$\psi^{-1}(b)$, $b \in \omega'$, are the blocks of treatments declared
indistinguishable.
\item $\psi$ \emph{injective and not bijective}: \emph{selection} of a proper
subfamily of treatments, i.e.\ the restriction of the experiment to some of
its groups. When $\psi$ is a proper inclusion $\omega' \subsetneq \omega$,
the surviving labels carried unchanged, we call it an \emph{insertion}
(design-reduction) morphism, after
Tjur~\cite[p.~1299]{tjur_discussion_2002}: his elementary criterion
quantifies over removals of sample points, and on the quantities of
\S\ref{sec:meaningful} that quantification acts through precisely this
class (Lemma~\ref{lem:tjurshadow}), a proper subclass of the morphisms over
which McCullagh's naturality quantifies; \S\ref{sec:meaningful} keeps the
two apart. Every injection is a bijection followed by an insertion, so
selections in general are relabellings composed with insertions.
\item $\psi$ \emph{bijective}: \emph{relabelling} of the treatments.
\item $\psi$ \emph{constant}: collapse of all treatments to a single one, the
design in which no treatment distinction is recorded.
\end{itemize}
Smaller choices of $\catO$ are legitimate and impose weaker naturality
demands: $\catO_{\mathrm{merge}} \subset \catO_{\mathrm{lab}}$, the wide
subcategory of surjections alone, and the two-object truncations
$\catO_{\mathrm{arr}} \subset \catO_{\mathrm{full}}$ used in
Example~\ref{ex:oneway}. Since naturality over a category restricts to
naturality over any subcategory, the verification below covers all of them at
once.

\item \textbf{Response scales.} $\catV$ has the single object $\R$ with its
Borel $\sigma$-algebra and only the identity morphism, so that the
response-scale naturality condition (displayed as~\eqref{eq:NV} in
\S\ref{sec:worked}) is vacuous. Affine equivariance of the one-way layout is true, but it is not an instance of \eqref{eq:NV}: see
Remark~\ref{cav:scaleV}.

\item \textbf{Sample spaces.} $\catS = \catV \times \catU^{\op}$ and
$\Gamma(\R,u) = \R^{u}$ with the product $\sigma$-algebra, the space of real
response vectors $y = (y_i)_{i \in u}$. On $\varphi^{\op} \colon u' \to u$ in
$\catU^{\op}$, $\Gamma$ gives restriction $\res_{\varphi} \colon \R^{u'} \to
\R^{u}$, $y \mapsto y \circ \varphi$.

\item \textbf{Designs.} $\cat{C} = \Set$ and both $J$'s are the
underlying-set functors, so an object of $\catD$ is a triple $(u,\omega,x)$
with $x \colon u \to \omega$ the \emph{treatment assignment}: $x(i)$ is the
treatment received by unit $i$. Equivalently, $x$ is the partition of $u$ into
\emph{groups}
\[
u_a \;\coloneqq\; x^{-1}(a) \subseteq u, \qquad a \in \omega,
\qquad n_a \coloneqq |u_a|, \qquad \textstyle\sum_{a \in \omega} n_a = n,
\]
indexed by $\omega$ and allowed to have empty parts. The design is
\emph{balanced} when all $n_a$ are equal, and $x$ is surjective exactly when
every treatment is observed. A morphism $(\varphi,\psi) \colon
(u,\omega,x) \to (u',\omega',x')$ is a subsample inclusion together with a
treatment map such that
\[
x'(\varphi(i)) \;=\; \psi(x(i)) \qquad (i \in u),
\]
the comma condition, displayed as~\eqref{eq:commacond} in \S\ref{sec:worked}:
\emph{each retained unit carries, in the larger design, the $\psi$-image of
the treatment it carried in the smaller one.}

\item \textbf{Parameters.} $\catK$ has objects
$W \times \R_{>0}$ with $W \in \Vect_{\R}^{\mathrm{fd}}$, morphisms
$A \times \id_{\R_{>0}}$ with $A$ linear, and $U_{\catK}$ the underlying set,
faithful (the same parameter category is used in Example~\ref{ex:linreg}
below). The parameter functor is
\[
\Theta(\omega) \;=\; \R^{\omega} \times \R_{>0},
\qquad
\Theta(\psi)(\mu',\sigma^{2}) \;=\; (\mu' \circ \psi,\ \sigma^{2}),
\quad\text{i.e.}\quad
\bigl(\Theta(\psi)\mu'\bigr)_{a} = \mu'_{\psi(a)} .
\]
A parameter value is a pair $\theta = (\mu,\sigma^{2})$ with $\mu =
(\mu_a)_{a \in \omega}$ the vector of \emph{group means}, one real number
per treatment label, not per unit, and $\sigma^{2}$ a single
\emph{common} (homoscedastic) error variance, shared by all groups and
therefore untouched by $\Theta(\psi)$. Precomposition $\psi^{*}$ is linear, so
$\Theta$ does land in $\catK$, and it is contravariant because a map of label
sets pulls means back, not forward.
\end{enumerate}

\noindent\textbf{The model.} For $d = (u,\omega,x)$ and $\theta =
(\mu,\sigma^{2}) \in \R^{\omega} \times \R_{>0}$,
\[
P_{\R,d}\bigl(\,\cdot \mid \mu,\sigma^{2}\bigr)
\;=\; \bigotimes_{i \in u} \Norm(\mu_{x(i)},\,\sigma^{2})
\qquad \text{on } \R^{u},
\]
that is, $Y_i = \mu_{x(i)} + \varepsilon_i$ with $\varepsilon_i$ independent
$\Norm(0,\sigma^{2})$: the textbook Gaussian one-way layout, with group sizes
$n_a$ read off the design and no constraint (such as $\sum_a \mu_a = 0$)
imposed on $\mu$.
\end{example}


\medskip
\noindent The example is carried through the whole paper: as a family of
Markov kernels in Example~\ref{ex:oneway-cont}, as the test case for
meaningfulness in \S\ref{sec:meaningful}, and, after truncating $\catO$ to two
objects and discretising $\Theta$, as the finite Baues--Wirsching computation
of Example~\ref{ex:oneway}.

\subsection{Br\o ns' comma-category equivalence}

Br\o ns' key structural observation is that the somewhat unwieldy
Definition~\ref{def:mcbmodel} is equivalent to a diagram in a single fixed
category of statistical models. Recall that for a functor
$F \colon \cat{B} \to \cat{E}$, the comma category $(\id_{\cat E} \downarrow
F)$ has objects the triples $(e, b, p)$ with $p \colon e \to F b$ in
$\cat{E}$.

\begin{definition}\label{def:statmod}
The \emph{category of statistical models} is
\[
\StatMod \;\coloneqq\; \bigl(\id_{\Set} \downarrow \Prb\,\bigr)
\]
formed from $\Prb\colon\Meas\to\Set$: an object is a set $\Theta_0$, a
measurable space $S$, and a map $\Theta_0 \to \Prb(S)$, i.e.\ exactly a
classical parametrised family of probability measures; a morphism is a pair
(reparametrisation, measurable map of sample spaces) making the square
commute.
\end{definition}

\begin{proposition}[after Br\o ns~{\cite[pp.~1280,~1282]{brons_discussion_2002}}]\label{prop:brons}
There is a bijection, natural in all the data, between
\begin{enumerate}
\item natural transformations $P \colon L \Rightarrow R$ as in
Definition~\ref{def:mcbmodel}, and
\item functors $\widehat{P} \colon \catV \times \catD^{\op} \to \StatMod$
whose composites with the two projections of the comma category are the given
functors $U_\catK \Theta \pr_\catO^{\op}\pr_2$ (to $\Set$) and
$\Gamma(\id \times \pr_\catU^{\op})$ (to $\Meas$).
\end{enumerate}
\end{proposition}

\begin{proof}
This is an instance of the universal property of the comma category: for
functors $G \colon \cat{A} \to \cat{E}$, $H \colon \cat{A} \to \cat{B}$ and
$F \colon \cat{B} \to \cat{E}$, lifts of $(G,H)$ to
$\cat{A} \to (\id_{\cat E} \downarrow F)$ commuting with both projections
correspond bijectively to natural transformations $G \Rightarrow F H$. Apply
this with $\cat{A} = \catV \times \catD^{\op}$, $\cat{E}=\Set$,
$\cat{B}=\Meas$, $G = L$, $H = \Gamma(\id\times\pr_\catU^{\op})$, $F = \Prb$: an
object assignment $a \mapsto (G a, H a, P_a)$ is functorial precisely when
each $P$-square commutes, which is naturality of $P$; morphism-level data of
$\widehat P$ are forced, so the correspondence is bijective. Naturality in the
data is a routine check.
\end{proof}

Br\o ns states the result with a sketch; the proof given above is complete.

\subsection{Two worked instantiations of the design data}
\label{sec:worked}

Definition~\ref{def:designdata} asks for six items, and it is easy to lose
sight of how ordinary the resulting objects are. We therefore instantiate it
twice, for multiple linear regression and for binary classification, verify
the naturality conditions of Definition~\ref{def:mcbmodel} in each case, and
record two non-examples. Throughout, $\FinSet_{\mathrm{inj}}$ denotes finite
sets and injections and $\Vect_{\R}^{\mathrm{fd}}$ finite-dimensional real
vector spaces and linear maps.

\begin{center}
\small
\begin{tabular}{@{}clll@{}}
\toprule
& item of Def.~\ref{def:designdata} & A: linear regression & B: binary classification \\
\midrule
(1) & $\catU$ & $\FinSet_{\mathrm{inj}}$ & $\FinSet_{\mathrm{inj}}$ \\
(2) & $\catO$ & $\Vect_{\R}^{\mathrm{fd}}$ & $\Vect_{\R}^{\mathrm{fd}}$ \\
(3) & $\catV$ & $\{\R\}$, identities only & $\{\mathbf{2},\mathbf{1}\}$, with
$!\colon\mathbf{2}\to\mathbf{1}$ \\
(4) & $\Gamma(v,u)$ & $\R^{u}$, Borel product & $v^{u}$, discrete \\
(5) & $\catD$ & $(J_{\catU} \downarrow J_{\catO})$: design matrices & idem: feature matrices \\
(6) & $\Theta(\omega)$ & $\omega^{*}\times\R_{>0}$ & $\omega^{*}$ \\
\bottomrule
\end{tabular}
\end{center}

In both cases $\cat{C} = \Set$ and $J_{\catU}$, $J_{\catO}$ are the
underlying-set functors, so an object of $\catD$ is a triple $(u,\omega,x)$
with $x \colon u \to \omega$ a set map, that is, exactly a design or
feature matrix with rows indexed by units, and a morphism
$(u,\omega,x)\to(u',\omega',x')$ is a pair $(\varphi,\psi)$ with
\begin{equation}\label{eq:commacond}
x' \circ \varphi \;=\; \psi \circ x .
\end{equation}
Condition~\eqref{eq:commacond}, together with the injectivity of $\varphi$
built into $\catU = \FinSet_{\mathrm{inj}}$ and the finiteness of $u$, is all
that the verifications below consume (injectivity is used at the
marginalisation step, and is necessary: for a non-injective $\varphi$
square~\eqref{eq:ND} fails in all three models, two units with the same image
becoming identical rather than independent). This is worth emphasising: the
comma-category definition
of $\catD$ is not bookkeeping; it is precisely the compatibility a model
needs.

Fix $v \in \catV$ and let
$f = (\varphi,\psi) \colon d=(u,\omega,x) \to d'=(u',\omega',x')$ in $\catD$.
In $\catD^{\op}$ this is an arrow $d' \to d$, and the square required by
Definition~\ref{def:mcbmodel} is
\begin{equation}\label{eq:ND}
\begin{tikzcd}[column sep=5em, row sep=2.6em]
U_{\catK}\Theta(\omega') \arrow[r,"P_{v,d'}"] \arrow[d,"\Theta(\psi)"'] &
\Prb\bigl(\Gamma(v,u')\bigr) \arrow[d,"\Prb(\res_{\varphi})"] \\
U_{\catK}\Theta(\omega) \arrow[r,"P_{v,d}"'] &
\Prb\bigl(\Gamma(v,u)\bigr)
\end{tikzcd}
\end{equation}
where $\res_{\varphi} \colon v^{u'} \to v^{u}$, $y \mapsto y\circ\varphi$, is
restriction along $\varphi$. In words: \emph{marginalise the big model onto
the subsample and one obtains the small model at the pulled-back parameter.}
Separately, for $\alpha \colon v \to v'$ in $\catV$ at fixed $d$, since $L$
does not depend on the $\catV$ coordinate, the square collapses to
\begin{equation}\label{eq:NV}
\Prb(\alpha^{u}) \circ P_{v,d} \;=\; P_{v',d}
\qquad\text{on } U_{\catK}\Theta(\omega).
\end{equation}

\subsubsection*{A. Multiple linear regression}

\begin{example}[Gaussian linear model]\label{ex:linreg}
\leavevmode
\begin{enumerate}
\item \textbf{Units.} $\catU = \FinSet_{\mathrm{inj}}$; an object $u$ indexes
a sample, a morphism $\varphi \colon u \hookrightarrow u'$ is the inclusion of
a subsample.

\item \textbf{Covariate spaces.} $\catO = \Vect_{\R}^{\mathrm{fd}}$. An object
$\omega \cong \R^{p}$ is the space in which covariate vectors live; a morphism
$\psi \colon \omega \to \omega'$ is the inclusion of a submodel's covariates
into a larger set, or a linear recoding (centring, orthogonalisation, contrast
coding).

\item \textbf{Response scales.} $\catV$ has the single object $\R$ (Borel) and
only the identity morphism; see Remark~\ref{cav:scaleV} for why the affine
group cannot be used here without amending Definition~\ref{def:designdata}.

\item \textbf{Sample spaces.} $\catS = \catV\times\catU^{\op}$ and
$\Gamma(\R,u) = \R^{u}$ with the product $\sigma$-algebra; on
$\varphi^{\op}\colon u' \to u$ in $\catU^{\op}$, $\Gamma$ gives restriction
$\res_\varphi \colon \R^{u'} \to \R^{u}$.

\item \textbf{Designs.} $\catD = (J_{\catU}\downarrow J_{\catO})$; an object
$(u,\omega,x)$ is a design matrix, $x(i) \in \omega$ being the covariate
vector of unit $i$. Choosing a basis of $\omega\cong\R^{p}$ and an enumeration
of $u$, the map $x$ \emph{is} the $n\times p$ matrix $X$.

\item \textbf{Parameters.} $\catK$ has objects $W\times\R_{>0}$ for
$W \in \Vect_{\R}^{\mathrm{fd}}$ and morphisms $A\times\id_{\R_{>0}}$ with $A$
linear; $U_{\catK}$ is the underlying set, faithful because a linear map is
determined by its underlying function. The parameter functor is the dual, the
subrepresentation of linear functionals inside McCullagh's standard
representation~\cite[\S4.2, p.~1235]{mccullagh_what_2002},
\[
\Theta(\omega) \;=\; \omega^{*}\times\R_{>0},
\qquad
\Theta(\psi)(\beta',\sigma^{2}) \;=\; (\beta'\circ\psi,\ \sigma^{2}) .
\]
\end{enumerate}

\noindent The model is, for $d=(u,\omega,x)$ and
$\theta = (\beta,\sigma^{2}) \in \omega^{*}\times\R_{>0}$,
\[
P_{\R,d}(\,\cdot \mid \beta,\sigma^{2})
\;=\; \bigotimes_{i\in u} \Norm(\beta(x(i)),\,\sigma^{2})
\qquad\text{on } \R^{u},
\]
that is, $Y \sim \Norm_{n}(X\beta,\sigma^{2}I_{n})$ in coordinates.
\end{example}

\begin{proof}[Verification of~\eqref{eq:ND}]
Take $f=(\varphi,\psi) \colon d \to d'$ and $\theta'=(\beta',\sigma^{2})$ on
$\omega'$. Down-then-across, $\Theta(\psi)\theta'
= (\beta'\circ\psi,\sigma^{2})$, so
\[
P_{\R,d}\bigl(\,\cdot \mid \Theta(\psi)\theta'\bigr)
= \bigotimes_{i\in u}\Norm(\beta'(\psi(x(i))),\sigma^{2}).
\]
Across-then-down, $P_{\R,d'}(\,\cdot\mid\theta')
= \bigotimes_{j\in u'}\Norm(\beta'(x'(j)),\sigma^{2})$, and pushing forward
along $\res_{\varphi}$ is marginalisation onto the coordinates
$\varphi(u) \subseteq u'$, which by independence gives
$\bigotimes_{i\in u}\Norm(\beta'(x'(\varphi(i))),\sigma^{2})$. The two agree
because $x'(\varphi(i)) = \psi(x(i))$, which is~\eqref{eq:commacond}. The
argument uses exactly two structural facts: independence across units, so that
marginalisation is computed coordinatewise, and the fact that the mean of unit
$i$ depends on the design only through $x(i)$. Condition~\eqref{eq:NV} is
vacuous because $\catV$ is discrete.
\end{proof}

\begin{remark}[Two familiar statistical statements, recovered]
Specialising~\eqref{eq:ND}: with $\psi = \id_{\omega}$ and $\varphi$ a proper
inclusion, the square says the model is consistent under subsampling, the
marginal of an $n$-observation regression on a subsample is the regression on
that subsample. With $\varphi = \id_{u}$ and
$\psi \colon \omega \hookrightarrow \omega'$ the inclusion of a submodel's
covariate space, it says that fitting the large model at a $\beta'$
annihilating the extra directions agrees with fitting the submodel at
$\beta'|_{\omega}$. Model nesting thus appears as a naturality condition
rather than as a constraint on coefficients.
\end{remark}

\begin{remark}[Intercept]
As stated, the model has no intercept. Two clean repairs: take $\catO$ to be
finite-dimensional real \emph{affine} spaces with affine maps and
$\Theta(\omega)$ the affine functionals on $\omega$ times $\R_{>0}$; or keep
$\Vect$ and restrict $\catD$ to designs whose covariate maps land in an affine
hyperplane $\{x_{0}=1\}$, with $\catO$ cut down to the linear maps preserving
it. The first is the more natural and is what the regression examples
of~\cite[\S6]{mccullagh_what_2002} effectively use.
\end{remark}

\begin{remark}[Design-conditional identifiability]\label{rem:identifiability}
$\Theta(\psi)=\psi^{*}\times\id$ is injective iff $\psi$ is surjective, and
$L(v,d)$ parametrises the design $d$ faithfully only relative to the image of
$x$. The design-conditional statement, that $\beta$ is identified from $d$
only on the span of $x(u) \subseteq \omega$, is invisible at the level of
Definition~\ref{def:designdata}; it is~\eqref{eq:ND} restricted to the
subcategory of designs of full column rank. This is the precise point at which
design-conditional sufficiency and identifiability enter the framework.
\end{remark}

\begin{caveat}[Response-scale morphisms express \emph{invariance}, not
equivariance]\label{cav:scaleV}
Because $L = U_{\catK}\Theta\,\pr_{\catO}^{\op}\pr_{2}$ has no $\catV$
dependence, a non-identity $\alpha \colon v \to v'$ imposes~\eqref{eq:NV},
namely $\Prb(\alpha^{u})\circ P_{v,d} = P_{v',d}$ \emph{at the same parameter
value}. For $\alpha(y)=ay+b$ the left-hand side is
$\bigotimes_{i}\Norm(a\beta(x(i))+b,\ a^{2}\sigma^{2})$, which is not
$P_{\R,d}(\cdot\mid\beta,\sigma^{2})$ unless $(a,b)=(1,0)$. The Gaussian
linear model is therefore \emph{not} natural over
$\catV = \mathrm{Aff}(\R)$, and taking $\catV$ discrete is the correct
instantiation.

Affine equivariance is of course a true and important property of the model,
but it is the statement that the model is a fixed point of an \emph{action} on
the parameter, which in the McCullagh--Br\o ns hexagon requires the parameter
functor to be defined on $\catV^{\op}\times\catO^{\op}$ rather than on
$\catO^{\op}$ alone. That is a mild generalisation of
Definition~\ref{def:designdata}, not an instance of it, and, as recorded in
the Remark following the hexagon in \S\ref{sec:mcb}, it is a return to
McCullagh's own construction rather than a departure from it: his parameter is
a subrepresentation of $v^{\omega}$ and so moves with the response scale,
whereas the $\catV$-independence of $\Theta$ is an artefact of Br\o ns'
packaging. In particular
Example~\ref{ex:oneway-setup} takes $\catV$
discrete for exactly this reason: the affine equivariance of the one-way
layout is a true statement about an action on $\Theta$, and lies outside
Definition~\ref{def:designdata} as written, whereas its equivariance under
\emph{relabelling of treatments} is genuine naturality, because that acts
through $\catO^{\op}$.
\end{caveat}

\subsubsection*{B. Binary classification}

\begin{example}[Bernoulli model with a linear predictor]\label{ex:logreg}
\leavevmode
\begin{enumerate}
\item \textbf{Units.} $\catU = \FinSet_{\mathrm{inj}}$, as above.

\item \textbf{Covariate spaces.} $\catO = \Vect_{\R}^{\mathrm{fd}}$; objects
are feature spaces, morphisms are linear feature maps.

\item \textbf{Response scales.} $\catV$ has two objects,
$\mathbf{2}=\{0,1\}$ and $\mathbf{1}=\{*\}$, and as morphisms the identities
together with the unique map $!\colon\mathbf{2}\to\mathbf{1}$. Statistically:
one may discard the label, not relabel it.

\item \textbf{Sample spaces.} $\Gamma(v,u)=v^{u}$ with the discrete
$\sigma$-algebra, so $\Gamma(\mathbf{2},u)=\{0,1\}^{u}$ is the label space and
$\Gamma(\mathbf{1},u)$ a one-point space; on $\varphi^{\op}$, restriction as
before.

\item \textbf{Designs.} $\catD = (J_{\catU}\downarrow J_{\catO})$; an object
$(u,\omega,x)$ is the feature matrix, the labels being absent. The labels are
the random object and hence live in $\Gamma$, not in $\catD$, which is the
categorical form of the statement that classification is a model conditional
on the features.

\item \textbf{Parameters.} $\catK = \Vect_{\R}^{\mathrm{fd}}$ with underlying
set $U_{\catK}$, and $\Theta(\omega)=\omega^{*}$,
$\Theta(\psi)(\beta') = \beta'\circ\psi$: the weight covector pulls back along
feature maps.
\end{enumerate}

\noindent With $\expit(t)=(1+e^{-t})^{-1}$ the model is
\[
P_{\mathbf{2},d}(\,\cdot\mid\beta)
\;=\; \bigotimes_{i\in u}\Bern\!\bigl(\expit(\beta(x(i)))\bigr)
\quad\text{on }\{0,1\}^{u},
\qquad
P_{\mathbf{1},d}(\,\cdot\mid\beta) = \delta \ \text{on } \{*\}^{u}.
\]
\end{example}

\begin{proof}[Verification]
For~\eqref{eq:ND} at $v=\mathbf{2}$, marginalising a product of Bernoullis
onto $\varphi(u)$ gives the first family below, and the pulled-back parameter
gives the second,
\[
\bigotimes_{i\in u}\Bern\bigl(\expit(\beta'(x'(\varphi(i))))\bigr),
\qquad
\bigotimes_{i\in u}\Bern\bigl(\expit((\beta'\circ\psi)(x(i)))\bigr),
\]
and these agree by~\eqref{eq:commacond}. At $v=\mathbf{1}$ both sides are the
unique measure on a one-point space. For~\eqref{eq:NV} with
$!\colon\mathbf{2}\to\mathbf{1}$, the pushforward of
$P_{\mathbf{2},d}(\cdot\mid\beta)$ along
$!^{u}\colon\{0,1\}^{u}\to\{*\}^{u}$ is the Dirac measure, which is
$P_{\mathbf{1},d}(\cdot\mid\beta)$ for every $\beta$.
\end{proof}

\begin{remark}[The link plays no role]\label{rem:glm}
The verification uses only that the inverse link is applied \emph{pointwise to
the linear predictor}: any measurable $g^{-1}\colon\R\to[0,1]$ yields a
natural family by the same argument. Naturality in $\catD$ is thus the
categorical content of the statement that a generalised linear model is
determined by a link applied unit by unit, and conversely that any
sample-level coupling of the predictions destroys it (Example~\ref{ex:softmax}).
\end{remark}

\begin{remark}[Why the label swap is excluded from $\catV$]
The involution $\tau(y)=1-y$ on $\mathbf{2}$ satisfies
$\Prb(\tau^{u})P_{\mathbf{2},d}(\cdot\mid\beta)
= \bigotimes_{i}\Bern(\expit(-\beta(x(i))))
= P_{\mathbf{2},d}(\cdot\mid -\beta)$, which is not
$P_{\mathbf{2},d}(\cdot\mid\beta)$. By~\eqref{eq:NV}, this violates naturality.
The correct reading is not that the model is defective, but that relabelling is
an equivariance, the phenomenon of Remark~\ref{cav:scaleV} in its smallest
possible instance, and hence a convenient test case for whether the
$\catV$-dependent parameter functor is worth adopting.
\end{remark}

\subsubsection*{Non-examples}

Both families below are well-defined and design-indexed; neither is a
McCullagh--Br\o ns model, and they fail in the same slot for the same reason.

\begin{example}[Empirical centring of the covariates]\label{ex:centring}
On the data of Example~\ref{ex:linreg} set
\[
\widetilde{P}_{\R,d}(\,\cdot\mid\beta,\sigma^{2})
= \bigotimes_{i\in u}\Norm(\beta(x(i)-\bar{x}_{d}),\ \sigma^{2}),
\qquad
\bar{x}_{d} = \frac{1}{|u|}\sum_{i\in u} x(i).
\]
This fails~\eqref{eq:ND} already for $\psi=\id$ and $\varphi$ a proper
inclusion, since $\bar{x}_{d} \neq \bar{x}_{d'}$: the mean assigned to a given
unit changes when other units are removed. The failure is in the $\catU$ slot,
and it is the mechanism behind Tjur's design-averaged non-example
(Example~\ref{ex:alphabetaxbar}; cf.\ Remark~\ref{rem:tjurscope}). Note that $\bar{x}$ is a perfectly good
statistic; what is illegitimate is its appearance in the \emph{definition} of
the family.
\end{example}

\begin{example}[Sample-normalised classifier]\label{ex:softmax}
On the data of Example~\ref{ex:logreg} let
\[
\widetilde{P}_{\mathbf{2},d}(\cdot\mid\beta)
= \bigotimes_{i\in u}\Bern\bigl(s_{i}(\beta,d)\bigr),
\qquad
s_{i} \;=\; \frac{e^{\beta(x(i))}}{\sum_{j\in u}e^{\beta(x(j))}},
\]
a softmax taken
across the sample rather than across the label set. Again the $\catU$ slot
fails: dropping a unit rescales every remaining probability. Normalisation
across a batch is a routine operation in learning pipelines, so this is not a
pathology invented for the purpose; its non-naturality is exactly the kind of
obstruction that the cohomological formulation of \S\ref{sec:hochschild}
detects and classifies.
\end{example}

\medskip
\noindent The three slots may now be summarised as follows.

\begin{center}
\small
\begin{tabular}{@{}lll@{}}
\toprule
slot & what naturality asserts & what it rules out \\
\midrule
$\catU$ & marginal onto a subsample is the subsample model
        & sample-dependent normalisation, empirical centring \\
$\catO$ & parameter pulls back along covariate maps
        & design-dependent reparametrisation \\
$\catV$ & \emph{invariance} at fixed parameter
        & relabelling, rescaling (these are equivariances) \\
\bottomrule
\end{tabular}
\end{center}

The verifications above are elementary. Remark~\ref{cav:scaleV} is the one
substantive point: with $\Theta$ defined on $\catO^{\op}$ alone,
response-scale morphisms can express only invariance, so textbook equivariance
statements about linear models lie outside Definition~\ref{def:designdata} as
written.

\begin{remark}[The one-way layout is Example~\ref{ex:linreg} in dummy
coding]\label{rem:dummy}
Write $\catD_{\mathrm{lab}}$ for the design category of
Example~\ref{ex:oneway-setup} and $\catD_{\mathrm{lin}}$ for that of
Example~\ref{ex:linreg}. Throughout this remark, $\omega$ denotes an object of
$\catO_{\mathrm{lab}}$, a finite treatment-label set as in
Example~\ref{ex:oneway-setup}, carrying no linear structure, \emph{not} the
covariate space $\omega \cong \R^{p}$ of Example~\ref{ex:linreg}, which
appears here as $\R[\omega]$.
Let $\R[-] \colon \FinSet \to \Vect_{\R}^{\mathrm{fd}}$ be the
free-vector-space functor and $\eta_{\omega} \colon \omega \to \R[\omega]$ the
basis inclusion, one-hot coding of a treatment label, which is natural,
$\R[\psi]\,\eta_{\omega} = \eta_{\omega'}\,\psi$. Then
\[
E(u,\omega,x) = \bigl(u,\ \R[\omega],\ \eta_{\omega} \circ x\bigr),
\qquad
E(\varphi,\psi) = (\varphi,\ \R[\psi])
\]
is a functor $\catD_{\mathrm{lab}} \to \catD_{\mathrm{lin}}$: it is
well defined on morphisms because $\R[\psi]\eta_{\omega}x =
\eta_{\omega'}\psi x = \eta_{\omega'}x'\varphi$, so~\eqref{eq:commacond} is
transported. Under the canonical isomorphism $\R[\omega]^{*} \cong
\R^{\omega}$, natural in $\omega$, one has $\Theta_{\mathrm{lab}} \cong
\Theta_{\mathrm{lin}} \circ \R[-]^{\op}$ and $\beta(\eta_{\omega}x(i)) = \mu_{x(i)}$, whence
\[
P^{\mathrm{one\text{-}way}}_{\R,\,d} \;=\; P^{\mathrm{lin}}_{\R,\,E d}
\qquad\text{for every } d \in \catD_{\mathrm{lab}} .
\]
The one-way layout is therefore the reindexing of the Gaussian linear model
along $E$, and its naturality is inherited, a natural transformation
restricting along any functor. Two things are \emph{not} inherited. First, the
converse fails, for two distinct reasons. $E$ is faithful but not full: a
morphism $(\varphi,\ell) \colon E(d) \to E(d')$ of $\catD_{\mathrm{lin}}$ must
send $\eta_\omega(x(i))$ to $\eta_{\omega'}(x'(\varphi(i)))$ for every unit
$i$, by~\eqref{eq:commacond}, which pins $\ell$ down to some $\R[\psi]$ on the
levels actually observed in $d$ but leaves it an arbitrary linear map on the
unobserved ones (when $x$ is surjective, every morphism from $E(d)$ to an
object in the image of $E$ is itself in the image of $E$). And
$\catD_{\mathrm{lin}}$ has objects outside the image of $E$ altogether,
recoded designs, whose covariate vectors are contrasts or sum-to-zero codings
rather than indicators, together with the morphisms into and among them
that no treatment map induces; centring, being affine rather than linear, is
not a morphism of $\catD_{\mathrm{lin}}$ at all. So
naturality over $\catD_{\mathrm{lin}}$ is strictly stronger than naturality
over $\catD_{\mathrm{lab}}$, and choosing label sets as $\catO$ is a
substantive weakening of the demand, not a change of notation. Second, the
cohomological test case of \S\ref{sec:hochschild} is computed over label
morphisms, so it certifies obstructions relative to that smaller class only;
this is one instance of the class-relativity insisted on in
\S\ref{sec:hochschild}.
\end{remark}

\medskip
\noindent\textbf{Example~\ref{ex:oneway-setup}, continued: verification of the
one-way layout.} With the naturality squares \eqref{eq:ND}--\eqref{eq:NV} now
displayed, we verify that the family of Example~\ref{ex:oneway-setup} is a
McCullagh--Br\o ns model, and record what each condition says for it.

\begin{proof}[Verification of~\eqref{eq:ND}]
Let $f = (\varphi,\psi) \colon d \to d'$ and $\theta' = (\mu',\sigma^{2})$ on
$\omega'$. Down-then-across, $\Theta(\psi)\theta' = (\mu' \circ
\psi,\sigma^{2})$, so
\[
P_{\R,d}\bigl(\,\cdot \mid \Theta(\psi)\theta'\bigr)
\;=\; \bigotimes_{i \in u} \Norm(\mu'_{\psi(x(i))},\ \sigma^{2}).
\]
Across-then-down, $P_{\R,d'}(\,\cdot \mid \theta') = \bigotimes_{j \in u'}
\Norm(\mu'_{x'(j)},\sigma^{2})$, and pushing forward along $\res_{\varphi}$ is
marginalisation onto the coordinates $\varphi(u) \subseteq u'$, which by
independence gives $\bigotimes_{i \in u}
\Norm(\mu'_{x'(\varphi(i))},\sigma^{2})$. The two agree because
$x'(\varphi(i)) = \psi(x(i))$, which is~\eqref{eq:commacond}. As in
Example~\ref{ex:linreg} the argument consumes exactly two structural facts:
independence across units, so that marginalisation is coordinatewise, and the
fact that the law of unit $i$ depends on the design only through $x(i)$.
Condition~\eqref{eq:NV} is vacuous because $\catV$ is discrete.
\end{proof}

\noindent\textbf{What naturality says, slot by slot.} Specialising the
naturality conditions~\eqref{eq:ND} and~\eqref{eq:NV} to morphisms living in
a single slot, and writing (N-$\cat{C}$) for the resulting condition on
morphisms of $\cat{C} \in \{\catU,\catO,\catV\}$:

\begin{itemize}
\item[(N-$\catU$)] $\psi = \id_{\omega}$, $\varphi$ a proper inclusion:
the marginal of the layout on the subsample $\varphi(u)$ is the layout on that
subsample, at the same $(\mu,\sigma^{2})$ and with the restricted assignment
$x' \circ \varphi$. \emph{Consistency under subsampling.}

\item[(N-$\catO$, merge)] $\varphi = \id_{u}$, $\psi$ surjective: the
pulled-back mean vector $\mu' \circ \psi$ is constant on the blocks of $\psi$,
and the square says that the fine-grained model evaluated at such a
block-constant mean vector \emph{is} the merged model at $\mu'$. Naturality
does not assert that merged groups have equal means; it asserts that on the
subfamily where they do, the two descriptions of the experiment coincide.

\item[(N-$\catO$, select)] $\psi$ injective: the same square read on a
subfamily of treatments. For $\psi$ an inclusion $\omega \subseteq \omega'$
this is design reduction, the covariate shadow of Tjur's sample-point
removals, and by Lemma~\ref{lem:tjurshadow} the only part of (N-$\catO$, select)
that his criterion reaches; a general injection adds a relabelling of the
surviving groups, which his criterion does not see. The distinction is not
idle: in Example~\ref{ex:oneway-findings}(i) the merge morphisms alone
certify nothing, and it is the selection morphisms that detect the failure.

\item[(N-$\catV$)] Vacuous here: by item~3 of
Example~\ref{ex:oneway-setup}, $\catV$ has only identity morphisms, so
\eqref{eq:NV} imposes nothing.
\end{itemize}

\noindent\textbf{Automorphisms: exchangeability and label equivariance.}
Taking $d' = d$ with $\varphi,\psi$ bijections, \eqref{eq:commacond} forces
$\varphi$ to carry $u_a$ onto $u_{\psi(a)}$, so
\[
\mathrm{Aut}_{\catD}(u,\omega,x)
\;=\;
\bigl\{(\varphi,\psi) : \psi \in \mathfrak{S}_{\omega},\ n_{\psi(a)} = n_a\ \forall a,\
\varphi|_{u_a} \colon u_a \xrightarrow{\ \sim\ } u_{\psi(a)}\bigr\},
\]
an extension of the group of size-preserving relabellings by
$\prod_{a \in \omega} \mathfrak{S}_{n_a}$, split once enumerations of the
groups are chosen. Naturality at $(\varphi,\id)$ is invariance of
$P_{\R,d}(\cdot \mid \theta)$ under permuting coordinates within groups, i.e.\
\emph{within-group exchangeability}; naturality at a general $(\varphi,\psi)$
is $(\res_{\varphi})_{*}P_{\R,d}(\cdot \mid \mu,\sigma^{2}) = P_{\R,d}(\cdot
\mid \mu \circ \psi,\sigma^{2})$, i.e.\ \emph{equivariance} under relabelling
of treatments. That equivariance is available in the $\catO$ slot but not
in the $\catV$ slot is not an accident of this example: $\Theta$ is defined on
$\catO^{\op}$, so a covariate morphism can act on the parameter, whereas a
response-scale morphism cannot (Remark~\ref{cav:scaleV}).

\noindent \textbf{Identifiability and design-conditional sufficiency.} $\Theta(\psi) = \psi^{*} \times \id$ is
injective iff $\psi$ is surjective and surjective iff $\psi$ is injective.
Independently of that, $\mu_a$ enters $P_{\R,d}$ only for $a \in x(u)$: at a
design that leaves some treatment unobserved, $P_{\R,d}$ factors through the
restriction $\R^{\omega} \to \R^{x(u)}$ and $\mu$ is identified only on the
observed levels, the exact analogue, for label sets, of the column-rank
condition of Remark~\ref{rem:identifiability}. The design-conditional
sufficient statistic is the pair (group sample means on $x(u)$, pooled
within-group sum of squares), and it is at this point that design-conditional
sufficiency and identifiability enter.

\section{Markov-categorical background}\label{sec:markov}

We use only a small fragment of the theory; see~\cite{fritz_synthetic_2020,perrone_markov_2024}
for the general framework.

The \emph{Giry monad} $(\G,\eta,\mu)$ is the probability-distribution monad on $\Meas$.

For a measurable space \(X=(X,\Sigma_X)\), $\G X$ is the set of all probability measures on \(X\). This set is itself equipped with a \(\sigma\)-algebra: namely, the smallest \(\sigma\)-algebra making all evaluation maps
\[
\operatorname{ev}_A:\mathcal G X\to [0,1],
\qquad
\operatorname{ev}_A(\mu)=\mu(A),
\]
measurable, for every \(A\in\Sigma_X\).

Thus \(\mathcal G X\) is again a measurable space.

If
\[
f: X\to Y
\]
is measurable, define
\[
\mathcal G f:\G X\to \G Y
\]
by pushforward:
\[
\G f(\pi)=f_\ast\pi,
\]
that is,
\[
(f_\ast\pi)(B)=\pi(f^{-1}(B)),
\qquad
B\in\Sigma_Y.
\]

Hence
\[
\mathcal G:\Meas\to\Meas
\]
is a functor.

The unit of the monad is
\[
\eta_X:X\to \G X,
\qquad
\eta_X(x)=\delta_x,
\]
where \(\delta_x\) denotes the Dirac probability measure concentrated at \(x\).

The multiplication of the monad is
\[
\mu_X:\G\G X\to\G X,
\]
defined by
\[
\mu_X(\Xi)(A)
=
\int_{\G X}\nu(A)\,d\Xi(\nu),
\qquad
A\in\Sigma_X.
\]

Thus an element \(\Xi\in\G\G X\), i.e.\ a probability distribution on the space of probability distributions on \(X\), is sent to its mixture.

Equivalently, for every bounded measurable function \(h:X\to\mathbb R\),
\[
\int_X h(x)\,d\mu_X(\Xi)(x)
=
\int_{\G X}
\left(
\int_X h(x)\,d\nu(x)
\right)
d\Xi(\nu).
\]

The triple
\[
(\G,\eta,\mu)
\]
satisfies the monad laws.

A Kleisli morphism for the Giry monad from \(X\) to \(Y\) is a measurable map
\[
K:X\to\G Y.
\]
Equivalently, it is a Markov kernel
\[
K:X\times\Sigma_Y\to[0,1]
\]
such that, for every \(x\in X\),
\[
B\mapsto K(x,B)
\]
is a probability measure on \(Y\), and for every \(B\in\Sigma_Y\),
\[
x\mapsto K(x,B)
\]
is measurable.

Given two Kleisli morphisms
\[
K:X\to\G Y,
\qquad
L:Y\to\G Z,
\]
their Kleisli composition is
\[
L\odot K:X\to\G Z
\]
defined by
\[
(L\odot K)(x)(C)
=
\int_Y L(y)(C)\,dK(x)(y),
\qquad
C\in\Sigma_Z.
\]

In other words, the Kleisli category of the Giry monad has measurable spaces as objects and Markov kernels as morphisms. This leads to the following definition.

\begin{definition}
The Kleisli category of the Giry monad is called $\Stoch$, and it is the category whose objects are measurable spaces and whose morphisms $f \colon X \rightsquigarrow Y$ are Markov kernels
$f \colon X \times \Sigma_Y \to [0,1]$, composed by the
Chapman--Kolmogorov formula, see~\cite{fritz_synthetic_2020, giry_categorical_1982,perrone_markov_2024}.
There is a natural bijection
\begin{equation}\label{eq:kleisli}
\Stoch(X, Y) \;\cong\; \Meas\bigl(X, \G Y\bigr).
\end{equation}

$\FinStoch$ is the full subcategory of finite discrete measurable spaces;
morphisms are stochastic matrices. 

$\Stoch$ is a Markov category: symmetric
monoidal under the product of measurable spaces, with copy and discard maps at
every object, naturality being imposed on discard
alone~\cite[Def.~2.1]{fritz_synthetic_2020}; we write
$\mathrm{copy}_{X} \colon X \to X \otimes X$ for the copy map at $X$.

A morphism $f \colon X \to Y$ of a Markov category is \emph{deterministic} if
it is a homomorphism of the copy comonoids, that is, if
\[
\mathrm{copy}_{Y} \circ f \;=\; (f \otimes f) \circ \mathrm{copy}_{X}
\]
\cite[Def.~10.1]{fritz_synthetic_2020}. 
Read left to right, the condition
says that running $f$ once and duplicating its output agrees with duplicating
the input and running $f$ twice independently; any genuine randomness in $f$
breaks this, since two independent runs need not agree. In $\Stoch$ the
condition says exactly that $f(x,B) \in \{0,1\}$ for every $x \in X$ and every
$B \in \Sigma_{Y}$, i.e.\ that each $f(x,\cdot)$ is a $\{0,1\}$-valued
measure~\cite[Ex.~10.4]{fritz_synthetic_2020}; in $\FinStoch$, that the
stochastic matrix of $f$ has entries in $\{0,1\}$. Deterministic morphisms
contain the identities and are closed under composition and under $\otimes$,
so they form a wide symmetric monoidal subcategory
$\Stoch_{\det} \subseteq \Stoch$.

\end{definition}


\begin{definition}\label{def:dirac}
The \emph{Dirac embedding} $\del \colon \Meas \to \Stoch$ is the identity on
objects and sends a measurable map $f$ to the kernel
$x \mapsto \delta_{f(x)}$. It is strong monoidal and lands in
$\Stoch_{\det}$, since a Dirac kernel is $\{0,1\}$-valued; it is faithful on
the measurable maps into $Y$ exactly when $\Sigma_{Y}$ separates points (so on
all standard Borel, discrete and finite codomains, but not on $\Meas$ as a
whole). The converse,
that every deterministic $X \rightsquigarrow Y$ is $\del$ of a measurable map,
is not automatic: it asks that the $\{0,1\}$-valued measures on $Y$ be
exactly its points. This holds when $\Sigma_{Y}$ is countably generated and
separates points, in particular for standard Borel $Y$, and always in
$\FinStoch$~\cite[Ex.~10.4]{fritz_synthetic_2020}. Hypotheses~(a) and~(b) of
Proposition~\ref{wprop:copower} are the two instances of this requirement that
arise below; see Remark~\ref{cav:Bscope} (ii).
\end{definition}

Two facts are used below without further comment: $(\mathrm{i})$
correspondence~\eqref{eq:kleisli} is natural in $X$ (via precomposition with
measurable maps) and in $Y$ (via pushforward); $(\mathrm{ii})$ a diagram of
Dirac kernels $\del f_{j}$ commutes in $\Stoch$ iff the diagram of the
$f_{j}$ commutes in $\Meas$, provided the codomains involved separate points,
since $\del$ is faithful.

\begin{remark}[What Markov categories add]\label{rem:whatmarkovadds}
The point of the change of codomain in Section~\ref{sec:bridge} is not
notational. Once a model lands in $\Stoch$, the synthetic apparatus
of~\cite{fritz_synthetic_2020}, a.s.-equality and the synthetic Fisher--Neyman
characterisation of sufficiency (which requires strict positivity, available in
$\Stoch$~\cite[Ex.~13.19, Thm.~14.5]{fritz_synthetic_2020}, not conditionals,
which $\Stoch$ lacks in general~\cite[Ex.~11.3]{fritz_synthetic_2020}), applies to each
component $P_{v,d}$ \emph{inside} the category, while McCullagh--Br\o ns
naturality constrains the family \emph{across} components. The two mechanisms
compose: e.g.\ sufficiency of a statistic for every design, plus naturality of
the statistic, is a statement with actual content about the whole diagram,
not a componentwise remark. Making such combined statements provable is the
purpose of the bridge.
\end{remark}

\section{The bridge}\label{sec:bridge}

The observation driving everything is: by~\eqref{eq:kleisli}, a
component $P_{v,d} \colon U_\catK\Theta(\omega) \to \Prb(\Gamma(v,u))$ of a
McCullagh--Br\o ns model \emph{wants} to be a morphism
$\Theta(\omega) \rightsquigarrow \Gamma(v,u)$ in $\Stoch$, the only gap between
Definition~\ref{def:mcbmodel} and that statement is measurability in
$\theta$, which the $\Set$-valued formulation cannot see. In the finite case
the gap vanishes, and the translation is an equivalence; in general it is an
added hypothesis, which we make explicit.

\subsection{Finite case}

Assume the data of Definition~\ref{def:designdata} are \emph{finite}: $\catV$,
$\catO$, $\catU$ take values with finite underlying sets, $\Gamma$ lands in finite
discrete spaces, and $U_\catK\Theta(\omega)$ is finite for every $\omega$.
Equip each $U_\catK\Theta(\omega)$ with its discrete measurable structure, so
that $U_\catK\Theta$ lifts (uniquely) to a functor
$\Theta_\Meas \colon \catO^{\op} \to \Meas$ with
$U_\catK \Theta = |\!-\!| \circ \Theta_\Meas,$ where $|\!-\!|:\Meas\to\Set$ is
the forgetful functor.

Define two functors into $\FinStoch$:
\[
\mathbf{L} \coloneqq \del \circ \Theta_\Meas \circ \pr_\catO^{\op} \circ \pr_2,
\qquad
\mathbf{R} \coloneqq \del \circ \Gamma \circ (\id_\catV \times \pr_\catU^{\op}),
\]
both $\catV \times \catD^{\op} \to \FinStoch$, where $\del$ is the Dirac
embedding of Definition~\ref{def:dirac}.

\begin{theorem}[Bridge, finite case]\label{thm:bridge-finite}
There is a canonical bijection between:
\begin{enumerate}
\item McCullagh--Br\o ns models $P \colon L \Rightarrow R$
(Definition~\ref{def:mcbmodel});
\item families of morphisms
$\mathbf{P}_{v,d} \colon \mathbf{L}(v,d) \rightsquigarrow \mathbf{R}(v,d)$ in
$\FinStoch$, indexed by the objects of $\catV\times\catD^{\op}$, such that for
every morphism $m$ of $\catV \times \catD^{\op}$ the square
\begin{equation}\label{eq:natsquare}
\begin{tikzcd}[column sep=3em]
\mathbf{L}(v,d) \arrow[r, rightsquigarrow, "\mathbf{P}_{v,d}"]
\arrow[d, "\mathbf{L}(m)"']
& \mathbf{R}(v,d) \arrow[d, "\mathbf{R}(m)"] \\
\mathbf{L}(v',d') \arrow[r, rightsquigarrow, "\mathbf{P}_{v',d'}"']
& \mathbf{R}(v',d')
\end{tikzcd}
\end{equation}
commutes in $\FinStoch$; i.e.\ natural transformations
$\mathbf{P}\colon \mathbf{L} \Rightarrow \mathbf{R}$.
\end{enumerate}
\end{theorem}

\begin{remark}[Where the probability sits in the square]\label{rem:squareshape}
In~(2), the vertical arrows $\mathbf{L}(m) = \del\,\Theta_\Meas(\psi)$ and
$\mathbf{R}(m) = \del\,\Gamma(\alpha,\varphi)$ are Dirac kernels, hence
deterministic, by the definition of $\mathbf{L}$ and $\mathbf{R}$; only the
components $\mathbf{P}_{v,d}$ carry probability, and these are arbitrary
stochastic matrices (Dirac only when the model is degenerate).
\end{remark}

\begin{proof}[Proof of Theorem~\ref{thm:bridge-finite}]
Write $d = (u,\omega,x)$, so that
$\mathbf{L}(v,d) = \del\Theta_\Meas(\omega)$ and
$\mathbf{R}(v,d) = \del\Gamma(v,u)$; since $\del$ is the identity on objects,
these are the measurable spaces $\Theta_\Meas(\omega)$ and $\Gamma(v,u)$,
both finite and discrete.

\emph{Step 1 (object level: transposition).} Fix $(v,d)$. We claim a bijection
\begin{equation}\label{eq:transpose}
\Set\bigl(U_\catK\Theta(\omega),\, \Prb(\Gamma(v,u))\bigr)
\;\xrightarrow{\ \cong\ }\;
\FinStoch\bigl(\Theta_\Meas(\omega),\, \Gamma(v,u)\bigr),
\qquad P_{v,d} \longmapsto \mathbf{P}_{v,d}.
\end{equation}
Indeed $U_\catK\Theta(\omega)$ is by construction the underlying set of the
discrete space $\Theta_\Meas(\omega)$, and $\Prb(S)$ is the underlying set of
the measurable space $\G S$. Because $\Theta_\Meas(\omega)$ is discrete,
\emph{every} function out of it is measurable, so
$\Set\bigl(U_\catK\Theta(\omega), \Prb(\Gamma(v,u))\bigr)
= \Meas\bigl(\Theta_\Meas(\omega), \G\Gamma(v,u)\bigr)$, an equality of
sets, not merely an inclusion. This is the component of the adjunction
$D \dashv U$ of Proposition~\ref{prop:coreflection} at
$\bigl(U_\catK\Theta(\omega), \G\Gamma(v,u)\bigr)$; what is done here by hand
at one pair of objects is done there uniformly. This is the only place where
discreteness of the parameter object is used, finiteness of the spaces
serves only to place $\mathbf{L}$, $\mathbf{R}$ and $\mathbf{P}$ in
$\FinStoch$ rather than in $\Stoch$, and it is exactly the gap
identified before Definition~\ref{def:mcbmodel}: in general one must
\emph{require} measurability in $\theta$. Composing with~\eqref{eq:kleisli}
gives~\eqref{eq:transpose}. Explicitly, $\mathbf{P}_{v,d}$ is the kernel
$\theta \mapsto P_{v,d}(\,\cdot \mid \theta)$, the same data read as a
morphism rather than as a $\Prb$-valued function.

\emph{Step 2 (morphism level: both sides are deterministic).} Let
$m = (\alpha, (\varphi,\psi)^{\op})$ be a morphism of
$\catV \times \catD^{\op}$, with target $(v',d')$, $d' = (u',\omega',x')$, so
that $(\varphi,\psi) \colon d' \to d$ in $\catD$, with
$\varphi \colon u' \to u$ and $\psi \colon \omega' \to \omega$. On
the $\Set$ side $L(m) = U_\catK\Theta(\psi)$ and
$R(m) = \Prb\bigl(\Gamma(\alpha,\varphi)\bigr)$, the latter being pushforward
of measures along $\Gamma(\alpha,\varphi)$. On the $\FinStoch$ side, by
definition of $\mathbf{L}$ and $\mathbf{R}$,
\[
\mathbf{L}(m) = \del\bigl(\Theta_\Meas(\psi)\bigr),
\qquad
\mathbf{R}(m) = \del\bigl(\Gamma(\alpha,\varphi)\bigr),
\]
so the vertical arrows of~\eqref{eq:natsquare} are Dirac kernels: $\mathbf{L}$
and $\mathbf{R}$ send every morphism to a deterministic kernel, while nothing
constrains $\mathbf{P}_{v,d}$ itself.

\emph{Step 3 (the two transposition rules).} Recall that composition in
$\Stoch$ is given by
$(k' \circ k)(x)(B) = \int k'(y)(B)\, k(x)(\mathrm{d}y)$.
Two special cases are all we need. For a measurable $g \colon Y \to Y'$ and a kernel
$k \colon X \rightsquigarrow Y$,
\begin{equation}\label{eq:push}
\bigl(\del(g) \circ k\bigr)(x)(B)
= \int \delta_{g(y)}(B)\, k(x)(\mathrm{d}y)
= k(x)\bigl(g^{-1}B\bigr)
= \bigl(\G(g) \circ k\bigr)(x)(B),
\end{equation}
so postcomposing with $\del(g)$ in $\FinStoch$ \emph{is} postcomposing with
the pushforward $\Prb(g) = \G(g)$ under~\eqref{eq:transpose}. Dually, for a
measurable $h \colon X' \to X$,
\begin{equation}\label{eq:pull}
\bigl(k \circ \del(h)\bigr)(x')(B)
= \int k(x)(B)\, \delta_{h(x')}(\mathrm{d}x)
= k\bigl(h(x')\bigr)(B),
\end{equation}
so precomposing with $\del(h)$ is precomposing with $h$ itself. These are
the two naturality statements of~\eqref{eq:kleisli} recorded after
Definition~\ref{def:dirac}, here written out in full.

\emph{Step 4 (squares correspond).} The $\Set$-naturality square for $P$ at
$m$ is the equality of functions
$U_\catK\Theta(\omega) \to \Prb(\Gamma(v',u'))$
\[
\Prb\bigl(\Gamma(\alpha,\varphi)\bigr) \circ P_{v,d}
\;=\; P_{v',d'} \circ U_\catK\Theta(\psi).
\]
Transpose each side under~\eqref{eq:transpose}. By~\eqref{eq:push} the
left-hand side transposes to
$\del\bigl(\Gamma(\alpha,\varphi)\bigr) \circ \mathbf{P}_{v,d}
= \mathbf{R}(m) \circ \mathbf{P}_{v,d}$; by~\eqref{eq:pull} the right-hand
side transposes to
$\mathbf{P}_{v',d'} \circ \del\bigl(\Theta_\Meas(\psi)\bigr)
= \mathbf{P}_{v',d'} \circ \mathbf{L}(m)$. Since~\eqref{eq:transpose} is a
bijection of hom-sets, the two functions are equal if and only if the two
kernels are equal, i.e.\ if and only if square~\eqref{eq:natsquare} commutes.

\emph{Step 5 (conclusion).} Steps~1 and~4 give, for each family, a family with
the same index set and the same commuting squares, and the assignment
$P_{v,d} \mapsto \mathbf{P}_{v,d}$ is invertible objectwise
by~\eqref{eq:transpose}, with the inverse governed by the same two rules read
right to left. Hence natural transformations $P \colon L \Rightarrow R$
correspond bijectively to natural transformations
$\mathbf{P} \colon \mathbf{L} \Rightarrow \mathbf{R}$. The final assertions
are Step~2 together with the observation that $\mathbf{P}_{v,d}$ lies in the
image of $\del$ precisely when each $P_{v,d}(\,\cdot \mid \theta)$ is a Dirac
measure, which is not required of a statistical model.
\end{proof}

The proof is a Kleisli-transposition argument; its value is not depth but the
licence it grants: every synthetic notion available in $\FinStoch$ applies to
finite McCullagh--Br\o ns models verbatim.



\medskip
\noindent\textbf{Is finiteness a conceptual restriction? Rounding as a
$\catV$-morphism.}
A natural objection to the finiteness hypothesis of
Theorem~\ref{thm:bridge-finite} is that it costs nothing: responses are
bounded by the range of any real instrument, covariate ranges are bounded for
the same reason, and a real number recorded in a $64$-bit floating-point
representation takes one of finitely many values, so the data of any actual
experiment are finite and Theorem~\ref{thm:bridge-finite} already covers
everything that is ever computed. Half of this is correct and worth recording
precisely; the other half is not, and the two halves are best kept apart.

\emph{What is correct.} Fix $M > 0$ and a finite grid $F \subset [-M,M]$,
for definiteness, the floating-point numbers representable in that range,
and let $\rho \colon \R \to F$ be rounding, with values outside $[-M,M]$
clamped to the extreme grid points. Then $\rho$ is Borel measurable and
surjective, hence precisely a coarsening morphism of the kind item~3 (the
response-scale category $\catV$) of Definition~\ref{def:designdata} admits.
Enlarge the response-scale category accordingly: let $\catV_{F}$ have the two
objects $\R$ and $F$, the grid carrying the discrete $\sigma$-algebra,
and, besides the identities, the single morphism $\rho \colon \R \to F$; this
is again a response-scale category in the sense of
Definition~\ref{def:designdata}(3), and $\Gamma$ extends canonically by
$\Gamma(F,u) = F^{u}$. Setting
\[
P_{F,d}(\,\cdot \mid \theta) \;\coloneqq\; (\rho^{u})_{*}\,P_{\R,d}(\,\cdot
\mid \theta)
\]
makes~\eqref{eq:NV} hold by construction, and every component now has finite
codomain. Two things follow. First, quantisation is not external to the
framework and not an approximation to the model: the discretised model
\emph{is} the image of the model under a morphism of the enlarged response-scale category $\catV_{F}$, obtained with no
hypothesis and no error term, and recorded data always live in the image of
such a morphism. Second, this is one of the rare non-identity response-scale morphisms
genuinely available to the Gaussian models of \S\ref{sec:worked}: unlike affine
rescaling, which acts on the parameter and therefore lies outside
Definition~\ref{def:designdata} as written (Remark~\ref{cav:scaleV}), rounding
is compatible with the model \emph{at the same parameter value}, which is
exactly what~\eqref{eq:NV} demands.

\emph{What is not correct.} The argument does not reduce
Theorem~\ref{thm:bridge-meas} to Theorem~\ref{thm:bridge-finite}, for three
independent reasons.

\emph{(a) It discretises the sample side only.} The hypothesis of
Theorem~\ref{thm:bridge-finite} also requires $U_\catK\Theta(\omega)$ to be
finite. Rounding is a morphism of $\catV$ and leaves the parameter functor
$\Theta$ untouched, whatever it is; in every model of \S\ref{sec:worked} the
parameter objects are continua, $\R^{\omega} \times \R_{>0}$ in the
one-way layout of Example~\ref{ex:oneway-setup}, for instance, so
$P_{F,d}$ is still a family indexed by a continuum: it is an object for
Theorem~\ref{thm:bridge-meas}, whose measurability hypothesis it inherits
from $P_{\R,d}$ rather than avoids.

\emph{(b) Discretising the parameter as well does not in general yield an
instance of Definition~\ref{def:designdata}.} Where $\Theta(\psi)$ is mere
reindexing, as in the one-way layout, a finite grid of parameter values is
preserved and Theorem~\ref{thm:bridge-finite} applies; this is the setting
of \S\ref{sec:hochschild}. Where $\Theta(\psi) = \psi^{*}$ uses the linear
structure of the covariate category, as in Examples~\ref{ex:linreg}
and~\ref{ex:logreg}, it does not: a floating-point grid is not closed under
arithmetic, addition is not associative on it and $\rho(a+b) \neq
\rho(a)+\rho(b)$ in general, so $F^{p}$ is not a vector space, the covariate
category $\Vect_{\R}^{\mathrm{fd}}$ of Example~\ref{ex:linreg} has no
floating-point analogue, and $\Theta(\psi) = \psi^{*}$ would be functorial only
up to rounding. Naturality is an \emph{equational} condition: the squares would
commute approximately, and no notion of approximate naturality is developed
here. The computation of \S\ref{sec:hochschild} is carried out in exact
arithmetic for this reason. Formulating an $\varepsilon$-naturality with a
quantitative obstruction is a reasonable question and would be the honest home
for numerical claims; nothing in this draft supplies one.

\emph{(c) The statistical content is not preserved by $\rho_{*}$.} The
pushforward of a Gaussian family along $\rho$ is not an exponential family, and
the pair (group means, within-group sum of squares) is not sufficient for it.
Statements about sufficiency and about exponential families are theorems about
the continuum object; its finite images do not inherit them.

\medskip
\emph{The correct relation is a limit, not a rounding.} Let $(\rho_n \colon S \to F_n)_{n}$ be a
refining sequence of finite quantisations whose cells generate the Borel
$\sigma$-algebra of $S$. Since the cells form a generating $\pi$-system, a
probability measure is determined by its pushforwards, so the canonical map
$\Prb(S) \to \varprojlim_n \Prb(F_n)$ is injective: nothing studied here is
invisible at finite level, and in this sense finiteness is indeed not a
conceptual limitation. But the map is not a bijection for free: the projective
limit $\hat S = \varprojlim_n F_n$ is a compact zero-dimensional space into
which $S$ embeds, and $\varprojlim_n \Prb(F_n) \cong \Prb(\hat S)$ by a
Kolmogorov-type extension theorem; $\Prb(S)$ is then the subset of measures
carried by the image of $S$, which is Borel for standard Borel $S$ but is
proper in general (for the dyadic quantisation of $[0,1)$ the consistent
family concentrating on the cells $[\tfrac12 - 2^{-n}, \tfrac12)$ has no
preimage). It is at exactly that point that the regularity hypothesis of
Theorem~\ref{thm:bridge-meas} reappears rather than disappears.
The measurable case is thus the projective limit of the finite ones, not a
special case of them, and cannot be dispensed with by quantising the data.

The injectivity of $\Prb(S) \to \varprojlim_n \Prb(F_n)$ is elementary; the
identification of the projective limit with $\Prb(\hat S)$ is classical but
is not used below. 

\subsection{Measurable case}

\begin{theorem}[Bridge, measurable case]\label{thm:bridge-meas}
Suppose the $\Set$-valued parameter functor
$U_\catK\Theta \colon \catO^{\op} \to \Set$ lifts to a functor
$\Theta_\Meas \colon \catO^{\op} \to \Meas$ along the underlying-set functor
$|\!-\!| \colon \Meas \to \Set$, that is,
$|\!-\!| \circ \Theta_\Meas = U_\catK \circ \Theta$ (i.e.\ each parameter
object carries a measurable structure and the reparametrisation maps
$\Theta(\psi)$ are measurable), and let
$P \colon L \Rightarrow R$ be a McCullagh--Br\o ns model such that each
component $P_{v,d} \colon U_\catK\Theta(\omega) \to \Prb(\Gamma(v,u))$ is
measurable as a map $\Theta_\Meas(\omega) \to \G \Gamma(v,u)$ (equivalently:
$\theta \mapsto P_{v,d}(A \mid \theta)$ is measurable for every measurable
$A$); call such a $P$ a \emph{measurable model}. Define
$\mathbf{L}, \mathbf{R} \colon \catV \times \catD^{\op} \to \Stoch$ by the
same formulas as in the finite case,
$\mathbf{L} = \del\,\Theta_\Meas\,\pr_\catO^{\op}\,\pr_2$ and
$\mathbf{R} = \del\,\Gamma\,(\id_\catV \times \pr_\catU^{\op})$. Then the
conclusion of Theorem~\ref{thm:bridge-finite} holds with $\FinStoch$ replaced
by $\Stoch$: there is a canonical bijection between measurable models and
natural transformations $\mathbf{P} \colon \mathbf{L} \Rightarrow \mathbf{R}$;
the arrows $\mathbf{L}(m)$, $\mathbf{R}(m)$ are Dirac kernels and hence
deterministic (Definition~\ref{def:dirac}), with no further hypothesis;
and each component $\mathbf{P}_{v,d}$ is a kernel that need not be
deterministic.
\end{theorem}

\begin{proof}
As in the finite case, with Step~1 replaced by the measurability hypothesis:
\eqref{eq:kleisli} and its naturality hold for arbitrary measurable spaces;
measurability of $P_{v,d}$ is exactly what places it in the hom-set
\[
\Meas\bigl(\Theta_\Meas(\omega),\,\G\Gamma(v,u)\bigr)
\]
that \eqref{eq:kleisli} transposes; and the square-transposition argument
(Steps~3--4) uses only naturality of the Kleisli bijection, not finiteness.
\end{proof}

\begin{remark}[The hypotheses are a genuine restriction]\label{rem:hypotheses}
McCullagh's formulation is $\Set$-valued: parameter objects are bare sets (or
objects of an abstract concrete category $\catK$), and no measurability in
$\theta$ is imposed. Theorem~\ref{thm:bridge-meas} therefore does \emph{not}
say ``every McCullagh--Br\o ns model lives in $\Stoch$''; it says the ones
satisfying a mild regularity hypothesis do. Two mitigations. (i) Every
statistical model used in practice, in particular every exponential-family
model, and all examples in~\cite{mccullagh_what_2002} with finite-dimensional
parameter spaces, satisfies the hypothesis with the obvious Borel structures
(the measure- and function-valued parameter spaces of~\cite[\S8]{mccullagh_what_2002}
would require a choice of $\sigma$-algebra, not pursued here); kernels $\theta \mapsto P_\theta$ with
measurable dependence are the standard regularity assumption of decision
theory. (ii) The hypothesis is not merely harmless but \emph{clarifying}: the
$\Set$-valued formulation silently permits models to which no measure-theoretic
statistical procedure (risk, Bayes, estimation) could be applied; the bridge
makes the regularity requirement part of the type. The bridge should
therefore be stated with this hypothesis, not without it.
\end{remark}

\begin{caveat}[No claim of essential surjectivity onto synthetic models]
The bridge exhibits McCullagh--Br\o ns models as certain diagrams in
$\Stoch$. It does not claim that every ``statistical model'' definable
synthetically in an abstract Markov category arises this way, nor that
$\Stoch$ is the only reasonable codomain (quasi-Borel spaces~\cite{heunen_convenient_2017}
or Borel spaces are alternatives with better categorical properties). The
choice of $\Stoch$ is the minimal one supporting the classic examples.
\end{caveat}

\begin{example}[One-way layout, continued]\label{ex:oneway-cont}
In Example~\ref{ex:oneway-setup}, $\Theta(\omega)=\R^\omega\times\R_{>0}$
with its Borel structure; measurability of
$(\mu,\sigma^2) \mapsto \bigotimes_i N(\mu_{x(i)},\sigma^2)$ is classical.
Theorem~\ref{thm:bridge-meas} applies, and the three naturality conditions
listed in Example~\ref{ex:oneway-setup} become commuting squares of kernels.
\end{example}


\subsection{The category of statistical models, reformulated}\label{sec:Bref}

It is natural to ask whether Br\o ns' category of statistical models is
identified with a category built canonically from $\Stoch$: objects its
morphisms, morphisms the commuting squares with deterministic sides. The following remark lets us say precisely
how much of it is true.

\begin{remark}[The bridge is a representability computation]\label{rem:repr}
The proof of Theorem~\ref{thm:bridge-finite} is elementary because it is, in disguise, the composite of
two standard facts. Isolating them is worth the space: they explain what the finiteness
hypothesis is doing, and they are what the rest of this section runs on.

\emph{(i) The probability functor is representable in $\Stoch$.} The monoidal
unit $I$ of $\Stoch$ is the one-point measurable space, and
\[
\Stoch(I,S) \;=\; \Meas(1,\G S) \;=\; \G S \;=\; \Prb(S),
\]
naturally in $S$; that is, $\Prb \cong \Stoch(I,-)$. A probability measure is a
generalised element at the unit, which is the precise content of the
standard slogan that probability measures are \emph{states}.

\emph{(ii) Discrete parameter sets are copowers of the unit.} The Dirac
embedding $\del$ (Definition~\ref{def:dirac}) is the free functor of the
Kleisli adjunction $\Meas \rightleftarrows \Stoch$, hence a left adjoint,
hence preserves colimits. Coproducts in $\Meas$ are disjoint unions (with the
disjoint-union $\sigma$-algebra), so for any set $A$,
\[
\del\bigl(A,2^A \bigr) \;\cong\; \coprod_{A} I \;=:\; A\cdot I .
\]

Combining, for any set $A$ and measurable space $S$,
\begin{equation}\label{eq:copower}
\Set\bigl(A,\Prb(S)\bigr)
\;\cong\;
\Set\bigl(A,\Stoch(I,S)\bigr)
\;\cong\;
\Stoch\bigl(A\cdot I,\;S\bigr),
\end{equation}
which is exactly the instance of~\eqref{eq:kleisli} used in the proof of
Theorem~\ref{thm:bridge-finite}.

Two consequences. First, determinism of the parameter side is not an extra
observation about the finite case: it is the statement that the domain object
is a copower of $I$ and that $\mathbf{L}$ sends every morphism to the Dirac
kernel of a map of the indexing sets, so $\mathbf{L}$ lands in the image of
$\del$ by construction.
Second, the finiteness hypothesis is doing two separable jobs, it makes
every set-map into $\Prb(S)$ a kernel (that is~\eqref{eq:copower}), and it
makes $\del$ full and faithful onto deterministic morphisms (that is the
converse recorded in Definition~\ref{def:dirac}), and only the second is a
genuine restriction,
since~\eqref{eq:copower} holds verbatim over arbitrary sets.

Proposition~\ref{prop:brons} admits the same reading one level up: the comma
category $(\id_{\Set}\downarrow\Prb)$ represents, as an object of $\Cat$ and
up to the evident $2$-categorical refinements, the
assignment $\cat{A} \mapsto \{(G,H,\alpha) \mid \alpha \colon G \Rightarrow
\Prb H\}$, and the universal property invoked in its proof is the
corresponding representability argument.
\end{remark}




\begin{definition}\label{def:arrdet}
Let $q \colon X \rightsquigarrow Y$ and $q' \colon X' \rightsquigarrow Y'$ be
morphisms of $\Stoch$. A \emph{square} from $q$ to $q'$ is a pair $(a,b)$ of
deterministic morphisms $a \colon X \rightsquigarrow X'$,
$b \colon Y \rightsquigarrow Y'$ such that
\[
\begin{tikzcd}[column sep=2.6em, row sep=2.2em]
X \arrow[r, squiggly, "q"] \arrow[d, squiggly, "a"'] & Y \arrow[d, squiggly, "b"] \\
X' \arrow[r, squiggly, "q'"'] & Y'
\end{tikzcd}
\qquad\qquad
q' \circ a \;=\; b \circ q \quad\text{in } \Stoch .
\]
We call $a$ and $b$ the \emph{legs} of the square. Squares compose
componentwise, $(a',b') \circ (a,b) = (a' \circ a,\, b' \circ b)$, and
$(\id_X,\id_Y)$ is a square from $q$ to itself; since the deterministic
morphisms form a subcategory $\Stoch_{\det}$ (\S\ref{sec:markov}), this
defines a category $\Arrdet{\Stoch}$ whose objects are \emph{all} the
morphisms of $\Stoch$ and whose morphisms are the squares. Two subcategories
are needed below:
\begin{enumerate}
\item[(i)] $\Arrcop{\Stoch}$, the full subcategory on the objects
$q \colon X \rightsquigarrow Y$ whose domain $X$ is isomorphic in $\Stoch$ to
a copower $A \cdot I$ of the monoidal unit;
\item[(ii)] $\Arrim{\Stoch}$, the wide subcategory whose morphisms are the
squares with both legs in the image of the Dirac embedding $\del$
(Definition~\ref{def:dirac}); it is closed under composition because $\del$
is a functor.
\end{enumerate}
\end{definition}

\begin{lemma}[Pointwise form of a square]\label{lem:square-pointwise}
Let $f \colon X \to X'$ and $g \colon Y \to Y'$ be measurable maps. Then
$(\del f, \del g)$ is a square from $q$ to $q'$ if and only if
\[
q'\bigl(f(x)\bigr) \;=\; g_{*}\, q(x)
\qquad\text{in } \Prb(Y'), \text{ for every } x \in X .
\]
\end{lemma}

\begin{proof}
By~\eqref{eq:pull}, $(q' \circ \del f)(x) = q'(f(x))$; by~\eqref{eq:push},
$(\del g \circ q)(x) = g_{*}\, q(x)$.
\end{proof}

\begin{remark}[Deterministic legs versus Dirac legs]\label{rem:arrim}
$\Arrim{\Stoch} \subseteq \Arrdet{\Stoch}$ because Dirac kernels are
deterministic, and the inclusion is an equality in $\FinStoch$. In general the
deterministic morphisms into a space $Z$ are all Dirac kernels exactly when
every $\{0,1\}$-valued measure on $Z$ is a Dirac measure, the converse
recorded in Definition~\ref{def:dirac}; if some $\{0,1\}$-valued measure
$\nu$ on $Z$ is not Dirac, the constant kernel at $\nu$ is deterministic but
not in the image of $\del$. Hence every square with legs into $X'$ and $Y'$
lies in $\Arrim{\Stoch}$ when this holds for $X'$ and for $Y'$, in particular
when they satisfy hypothesis~(b) of Proposition~\ref{wprop:copower}.
Hypothesis~(b) is sufficient but not necessary: hypothesis~(a) of the same
proposition describes discrete spaces on which every $\{0,1\}$-valued measure
is Dirac although the $\sigma$-algebra need not be countably generated. The
parameter set of the one-way layout is a case in point: $\R^{\omega} \times
\R_{>0}$ with the discrete $\sigma$-algebra, which is what $\Theta_0 \cdot I$
is for $\Theta_0 = U_\catK\Theta(\omega)$ of Example~\ref{ex:oneway-setup},
has the cardinality of the continuum, below the first measurable cardinal, so
every countably additive $\{0,1\}$-valued measure on it is Dirac; but a
countably generated $\sigma$-algebra has at most continuum many elements,
while the discrete one has $2^{|\R|}$.
\end{remark}

\begin{proposition}[Object-level form of the identification]\label{wprop:copower}
Transposition across~\eqref{eq:copower} extends to a functor
\[
\Jc \colon \StatMod \longrightarrow \Arrdet{\Stoch},
\qquad
(\Theta_0, S, p) \longmapsto
\bigl(\widehat{p} \colon \Theta_0 \cdot I \rightsquigarrow S\bigr),
\]
with the following properties.
\begin{enumerate}
\item $\Jc$ is injective on objects, and restricts to a bijection between the
objects of $\StatMod$ and those arrows of $\Stoch$ whose domain \emph{is} a
copower $A \cdot I$ of the monoidal unit; consequently $\Jc$ is essentially
surjective onto $\Arrcop{\Stoch}$, whose objects are the arrows with domain
merely isomorphic to such a copower (Definition~\ref{def:arrdet}).
\item $\Jc$ is bijective on hom-sets: for objects $x = (\Theta_0,S,p)$ and
$x' = (\Theta_0',S',p')$, the map $\StatMod(x,x') \to \Arrdet{\Stoch}(\Jc x,\Jc x')$
induced by $\Jc$ is a bijection provided
\begin{enumerate}
\item[(a)] every countably additive $\{0,1\}$-valued measure on
$(\Theta_0',\,2^{\Theta_0'})$ is a Dirac measure, automatic when
$|\Theta_0'|$ is smaller than the first measurable cardinal
(Lemma~\ref{lem:app-ulam}), hence always in the finite and countable cases;
and
\item[(b)] for a measurable space $Z$: $\Sigma_{Z}$ is countably generated and
separates points, so that for every $W$ the deterministic morphisms
$W \rightsquigarrow Z$ are exactly the Dirac kernels of measurable
maps~\cite[Ex.~10.4]{fritz_synthetic_2020}. We refer to this as ``(b) for
$Z$''; what is required here is (b) for $S'$.
\end{enumerate}
\item Let $\StatMod_{(a),(b)} \subseteq \StatMod$ be the full subcategory on
the objects $(\Theta_0,S,p)$ with $\Theta_0$ satisfying~(a) and $S$
satisfying~(b). Then $\Jc$ restricts to an equivalence between
$\StatMod_{(a),(b)}$ and the full subcategory of $\Arrcop{\Stoch}$ on the
arrows $X \rightsquigarrow S$ with $S$ satisfying~(b) and $X \cong A\cdot I$
for $A$ satisfying~(a). When all data are finite, write
$\StatMod_{\mathrm{fin}}$ for the full subcategory on finite $\Theta_0$ and
finite discrete $S$, (a) and~(b) are automatic, every object of
$\FinStoch$ is a copower of $I$, and
$\StatMod_{\mathrm{fin}} \simeq \Arrdet{\FinStoch}$ unconditionally.
\end{enumerate}
\end{proposition}

\begin{proof}[Proof idea]
With $A \cdot I \coloneqq (A, 2^{A})$, transposition~\eqref{eq:copower} gives
the object bijection of part~(1); passing to
$\Arrcop{\Stoch}$ requires knowing that an isomorphism
$A \cdot I \cong X$ of $\Stoch$ is deterministic, which is proved from the
copower side, where singletons are measurable. On morphisms,
$(r,f) \mapsto (\del r, \del f)$, the square commuting by
Lemma~\ref{lem:square-pointwise}; hypotheses~(a) and~(b) are exactly what is
needed to run that lemma backwards, turning the two deterministic legs of a
square into a function and a measurable map. The complete proof, together
with the facts about deterministic morphisms and $\{0,1\}$-valued measures it
relies on, is given in Appendix~\ref{app:copower}.
\end{proof}

The finite case is routine. The appeal to Ulam's theorem in~(a) should be
carried as an
explicit set-theoretic hypothesis, or avoided by restricting parameter objects
to standard Borel spaces throughout, which incurs no statistical cost.

\begin{remark}[Why the legs are deterministic]\label{rem:whydet}
Definition~\ref{def:arrdet} is not the most obvious choice. Write
$\Stoch^{\rightarrow}$ for the full arrow category, whose morphisms
$q \to q'$ are \emph{arbitrary} pairs of kernels $(a,b)$ with
$q' \circ a = b \circ q$; this is the definition one writes down first, and
$\Arrdet{\Stoch}$ is a wide, non-full subcategory of it. Three points on why
the restriction is imposed here.

\emph{(i) It is an image condition, inherited from the domain.} The morphisms
of $\StatMod$ are pairs of functions, and those of its measurable analogue
$\Statk$, introduced in Remark~\ref{cav:Bscope}(ii), pairs of measurable
maps, and transposition carries both to pairs of Dirac kernels. The image
of $\Jc$ therefore lies in $\Arrim{\Stoch}$, which coincides with
$\Arrdet{\Stoch}$ under the conditions recorded in
Remark~\ref{rem:arrim}, hypothesis~(b) of Proposition~\ref{wprop:copower}
among them; this is the precise sense in which the legs of
Definition~\ref{def:arrdet} are the transposed maps. Admitting arbitrary
pairs would make $\Jc$
non-full for trivial reasons and leave part~(2) of that proposition without
content. More structurally: Br\o ns' category is a comma category over a base
equipped with a probability functor, and the legs of its squares are morphisms
of that base; determinism is what survives of the base under transposition.

\emph{(ii) Two mechanisms are lost.} The pointwise form of the commutation
condition (Lemma~\ref{lem:square-pointwise}) rests on~\eqref{eq:push}
and~\eqref{eq:pull}, hence on the legs being Dirac; for stochastic legs there is no corresponding
condition in $\Meas$, and with it goes the fact that a diagram of
deterministic morphisms commutes in $\Stoch$ if and only if the underlying
diagram commutes in $\Meas$ (\S\ref{sec:markov}, fact~$(\mathrm{ii})$), the
mechanism by which $\Arrdet{\Stoch}$ is a faithful shadow of the
$\Meas$-level picture. Second, strict Kleisli equality, which is the right
condition for deterministic legs, is not obviously the right one once the legs
randomise: the ambient notion of equality in categorical probability is
almost-sure equality relative to a state~\cite[\S13]{fritz_synthetic_2020}, and
adopting it would make $\Stoch^{\rightarrow}$ a bicategory and any
identification a biequivalence. Neither obstacle is fatal, and both lie outside what is needed here.

\emph{(iii) The stochastic variant is a different subject, namely Le Cam's.}
Discussing~\cite{mccullagh_what_2002}, Huber recalls that handling sufficiency
forced Le Cam to replace pointwise maps by randomised maps, and that the
analysis of variance forced randomisation of the parameter spaces as well, so
that morphisms of experiments became equivalence classes of pairs of
transition functions~\cite{huber_discussion_2002}. His example separates the
two categories. Let $A$ be the experiment observing
$X = (X_1,\dots,X_n)$, $n \ge 2$, i.i.d.\ $\Norm(\mu,\sigma^{2})$, and $B$
the experiment observing independent $(U,V)$ with
$U \sim \Norm(\mu,\sigma^{2}/n)$ and $V/\sigma^{2} \sim \chi^{2}_{n-1}$: two
models with parameter set $\R \times \R_{>0}$, on $\R^{n}$ and on
$\R \times \R_{>0}$ respectively. The statistic $t = (\bar X, S^{2})$,
$S^{2} = \sum_i (X_i - \bar X)^{2}$, pushes each $P_A(\mu,\sigma^{2})$ to
$P_B(\mu,\sigma^{2})$, so $(\id, \del t)$ is a morphism $A \to B$ in
$\Arrdet{\Stoch}$. There is no morphism $B \to A$ in $\Arrdet{\Stoch}$,
whatever its parameter leg: if $g \colon \R \times \R_{>0} \to \R^{n}$ were
measurable with $g_{*}P_B(\theta) = P_A(a(\theta))$ for all $\theta$, the
direction $(g - \bar g\,(1,\dots,1))/S_g$ of the sample $g(U,V)$ would be
a statistic of $B$ whose law, uniform on the unit sphere of the hyperplane
$\sum_i x_i = 0$, is free of $\theta$; by Basu's
theorem~\cite{basu_statistics_1955} it would be independent of the complete
sufficient statistic $(U,V)$, of which it is a function, hence almost surely
constant, which it is not. In $\Stoch^{\rightarrow}$, by contrast, the
conditional distribution $k$ of $X$ given $(\bar X, S^{2})$, free of the
parameter, gives a morphism $(\id, k) \colon B \to A$, with
$t \circ k = \id_B$ but $k \circ t \neq \id_A$: $k \circ t$ sends
$x$ to the uniform distribution on the sphere through $x$. The two become
mutually inverse only after identifying morphisms that agree once composed
with the model, $k \circ t \circ p_A = p_A$, Huber's ``suitably defined
equivalence classes''. Allowing stochastic legs
thus coarsens the isomorphism relation on objects to Blackwell equivalence:
two models $p, p' \colon \Theta \rightsquigarrow \cdot$ on the same parameter
set are \emph{Blackwell equivalent} when each is a garbling of the other,
$p' = c \circ p$ and $p = c' \circ p'$ for kernels $c, c'$, Blackwell's
``equally informative''~\cite{blackwell_equivalent_1953}, deficiency zero in
Le Cam's theory~\cite{lecam_sufficiency_1964}, and mutual dominance for the
informativeness preorder of~\cite[\S5.1]{fritz_representable_2023}. In the
quotient just described, a morphism $B \to A$ over $\id_\Theta$ with an
inverse is exactly such a pair of garblings, so isomorphism there is
Blackwell equivalence, and the invariant under study becomes informativeness
rather than meaningfulness. That is the wrong invariant for the present purpose:
Definition~\ref{def:meaningful} asserts that reducing a design or
reparametrising introduces no randomness, and a category whose reductions may
themselves randomise cannot state it. The two readings of McCullagh's question,
which models are meaningful, and which experiments are equivalent, are
genuinely different questions, and this paper is committed to the first.
\end{remark}

\begin{caveat}[What this settles, and what it does not]\label{cav:Bscope}
Three points, in increasing order of consequence.

\emph{(i) The essential image is not everything.} $\StatMod$ is \emph{not}
equivalent to $\Arrdet{\Stoch}$: the image is the full subcategory on arrows
out of copowers of $I$. Any statement of the form ``Br\o ns' category
is the category whose objects are the morphisms of $\Stoch$'' is therefore
false as stated and must carry the domain restriction. The restriction is
invisible in the finite case only because there every object of $\FinStoch$ is
a copower of $I$.

\emph{(ii) The restriction is an artefact of the $\Set$-valued formulation.}
Replace Br\o ns' $\StatMod = (\id_{\Set}\downarrow\Prb)$ by its measurable
analogue $\Statk \coloneqq (\id_{\Meas}\downarrow\G)$, whose objects are
measurable maps $\Theta \to \G S$. By~\eqref{eq:kleisli} these are exactly the
morphisms $\Theta \rightsquigarrow S$ of $\Stoch$, with no copower condition,
and the same argument gives a functor
$\Jcm \colon \Statk \to \Arrdet{\Stoch}$
(Definition~\ref{def:arrdet}, whose objects are arbitrary, not
deterministic, morphisms of $\Stoch$). The two levels behave differently
and it is worth separating them.

On objects, the identification is a tautology, by~\eqref{eq:kleisli}: the
copower restriction disappears because every measurable space, and not only a
discrete one, is now admissible as a parameter object.

On morphisms it is not a tautology. Transposition sends a morphism $(u,h)$ of $\Statk$, a pair of
\emph{measurable} maps, to the square $(\del u, \del h)$, and this
assignment is unconditionally bijective on objects and full onto
$\Arrim{\Stoch}$. It is \emph{faithful} exactly when $\del$ is injective on
the measurable maps concerned, that is when $\Sigma_{\Theta'}$ and
$\Sigma_{S'}$ separate points; and $\Arrim{\Stoch} = \Arrdet{\Stoch}$ exactly
when every deterministic morphism into $\Theta'$ or into $S'$ is a Dirac
kernel (Remark~\ref{rem:arrim}), which is the remaining content of~(b). Both
requirements concern
\emph{both} legs, whereas Proposition~\ref{wprop:copower} invokes (b) for $S'$
alone, hypothesis~(a) doing the parameter-side work there, at the price of
the copower restriction. The correct statement is therefore: $\Jcm$
restricts to an equivalence between the full subcategory of $\Statk$ on the
objects whose parameter space and sample space both satisfy~(b) and the full
subcategory of $\Arrdet{\Stoch}$ on the arrows between such spaces;
unconditionally one has only a bijective-on-objects full functor
$\Statk \to \Arrim{\Stoch}$.

Hypothesis~(a) does not survive this passage, but neither does it drop
out: it served to recognise a deterministic morphism out of a \emph{discrete}
parameter object as a Dirac measure, and in the measurable setting that task
is taken over by (b) for $\Theta'$. The two are in fact instances of a single
requirement, the one isolated in Definition~\ref{def:dirac}: that the
$\{0,1\}$-valued measures on the codomain be exactly its points. Condition~(b)
secures it when the $\sigma$-algebra is countably generated and separating.
Condition~(a) is what remains of it over a discrete parameter object
$(\Theta_{0}', 2^{\Theta_{0}'})$, whose $\sigma$-algebra separates points but
is not countably generated unless $\Theta_{0}'$ is countable, so that the
identification of $\{0,1\}$-valued measures with points degenerates into a
set-theoretic question about measurable cardinals rather than a
measure-theoretic one. Passing to $\Statk$ therefore does not remove a
hypothesis: it replaces the set-theoretic instance by the measure-theoretic
one, and is repaid by the disappearance of the copower restriction. No copower
intervenes on either side. Nor
does any strengthening of the hypotheses give an equivalence with the full
arrow category $\Stoch^{\rightarrow}$: the morphisms of a comma category over
$\Meas$ are pairs of measurable maps, so their transposes are deterministic by
construction, and the genuinely stochastic squares are never in the image
(Remark~\ref{rem:whydet}).

So the unrestricted identification is correct \emph{for the measurable comma
category} on the full subcategory of objects satisfying~(b) on both legs, and
incorrect for the $\Set$-valued one, where it acquires a domain restriction.
The
mathematical content of the identification is thus not the equivalence itself but the
comparison functor $\StatMod \to \Statk$, which is precisely the
measurability gap of Remark~\ref{rem:hypotheses}, met from the other side.
That functor is constructed in Proposition~\ref{prop:coreflection} below. It
exists unconditionally, is fully faithful, and fails to be essentially
surjective exactly at the non-discrete parameter objects; so the copower
restriction of~(i) is the non-invertibility of a
counit, and it is invisible in the finite case for that reason.

\emph{(iii) This is a statement about $\StatMod$, not about models.} A
McCullagh--Br\o ns model is a diagram
$\widehat{P} \colon \catV\times\catD^{\op} \to \StatMod$
(Proposition~\ref{prop:brons}), and the parameter object at $\omega$ supplied
by Theorem~\ref{thm:bridge-finite} is $\del\Theta_\Meas(\omega)$, a copower of
$I$ exactly when $\Theta_\Meas(\omega)$ is discrete. The two identifications
therefore agree in the finite case and diverge in general.
Corollary~\ref{cor:coreflection-diagrams} says how far the diagram-level
statement can be pushed, and Proposition~\ref{wprop:copower} should not be
presented as going further than that.
\end{caveat}

The comparison functor of~(ii) can be written down, and doing so settles more
than the passage between the two formulations. Write
$D \colon \Set \to \Meas$ for the functor equipping a set with its power-set
$\sigma$-algebra and $U \colon \Meas \to \Set$ for the underlying-set functor.

\begin{proposition}[Discretisation of the parameter;
$\StatMod$ as a coreflective subcategory of $\Statk$]\label{prop:coreflection}
$D \dashv U$, and $U\G = \Prb$ by construction. Write
$\varepsilon \colon DU \Rightarrow \id_{\Meas}$ for the counit, the
component $\varepsilon_X$ is the identity function on $UX$, measurable because
its domain is discrete, and put
$\lambda \coloneqq \varepsilon\G \colon D\Prb \Rightarrow \G$. Then:
\begin{enumerate}
\item
$\Phi \colon \StatMod \to \Statk$,
$\;(\Theta_0, S, p) \mapsto (D\Theta_0,\, S,\, \lambda_S \circ Dp)$,
$\;(u,h) \mapsto (Du, h)$,
is a functor; its structure map is $p$ itself, read as a measurable map on a
discrete domain.
\item
$\Psi \colon \Statk \to \StatMod$,
$\;(\Theta, S, q) \mapsto (U\Theta,\, S,\, Uq)$,
$\;(a,h) \mapsto (Ua, h)$,
is a functor, and $\Phi \dashv \Psi$.
\item $\Psi\Phi = \id_{\StatMod}$. Hence the unit is an identity,
$\Phi$ is fully faithful, and $\StatMod$ is a coreflective subcategory of
$\Statk$ with coreflector $\Psi$.
\item The counit of $\Phi \dashv \Psi$ at $(\Theta, S, q)$ is
$(\varepsilon_{\Theta}, \id_S)$, invertible if and only if
$\Sigma_{\Theta} = 2^{\Theta}$. The essential image of $\Phi$ is the full
subcategory of $\Statk$ on the objects whose parameter space is discrete.
\end{enumerate}
None of this uses hypothesis~(a) or~(b) of Proposition~\ref{wprop:copower}, or
any set-theoretic assumption.
\end{proposition}

\begin{proof}
$D \dashv U$ because every function out of a discrete space is measurable, so
$\Meas(DA, X) = \Set(A, UX)$ naturally; $U\G = \Prb$ because the underlying
set of $\G S$ is $\Prb(S)$ and the underlying function of $\G(h)$ is
pushforward. As a whiskering of $\varepsilon$, $\lambda$ is natural, and each
$\lambda_S$ is the identity on underlying sets.

(1) $Dp \colon D\Theta_0 \to D\Prb(S)$ is measurable (both spaces discrete)
and $\lambda_S$ is measurable, so $\lambda_S \circ Dp$ is an object of
$\Statk$. For a morphism $(u,h)$, using naturality of $\lambda$, functoriality
of $D$, and the $\StatMod$ square $\Prb(h)\circ p = p' \circ u$,
\[
\G(h) \circ \lambda_S \circ Dp
= \lambda_{S'} \circ D\bigl(\Prb(h)\bigr) \circ Dp
= \lambda_{S'} \circ D\bigl(p' \circ u\bigr)
= \bigl(\lambda_{S'} \circ Dp'\bigr) \circ Du ,
\]
which is the $\Statk$ square. Identities and composites are preserved
componentwise.

(2) $\Psi$ is well defined since $Uq \colon U\Theta \to U\G S = \Prb(S)$. For
the adjunction, a morphism
$\Phi(\Theta_0,S,p) \to (\Theta',T,q)$ is a pair $(a,h)$ with $a \colon
D\Theta_0 \to \Theta'$ measurable and $q \circ a = \G(h) \circ \lambda_S \circ
Dp$. Transposing $a$ across $D \dashv U$ gives $u = Ua \colon \Theta_0 \to
U\Theta'$, and since both sides of the square are measurable maps out of the
discrete space $D\Theta_0$ and $U$ is faithful, the square is equivalent to
its image under $U$, namely $Uq \circ u = \Prb(h) \circ p$, the $\StatMod$
square for $(u,h) \colon (\Theta_0,S,p) \to \Psi(\Theta',T,q)$. The
correspondence is bijective, and natural in both variables because on the
parameter leg it is the $D \dashv U$ bijection and on the sample leg it is the
identity.

(3) $UD = \id_{\Set}$ and $U\lambda_S = \id$, so $\Psi\Phi$ is the identity on
objects and on morphisms. A left adjoint whose unit is invertible is fully
faithful, and its domain is then coreflective in its codomain.

(4) $\Phi\Psi(\Theta,S,q) = (DU\Theta,\, S,\, \lambda_S \circ D(Uq))$, and
$(\varepsilon_\Theta, \id_S)$ is a morphism from it to $(\Theta,S,q)$: both
legs of the required square are measurable maps $DU\Theta \to \G S$ with
underlying function $Uq$, and $U$ is faithful. It is the counit by the triangle
identities. Since $\varepsilon_\Theta$ is a bijection, it is an isomorphism in
$\Meas$ exactly when $\Sigma_\Theta$ is the full power set; and
$X \cong DA$ in $\Meas$ forces $\Sigma_X = 2^X$, which gives the description of
the essential image. Fullness of that subcategory is automatic from~(3).
\end{proof}

The point of~(4) is that the domain restriction which had to be carried
through Proposition~\ref{wprop:copower} by hand is a single counit, and one
whose failure to be invertible is exactly non-discreteness of the parameter.
Transporting to diagrams is then formal.

\begin{corollary}[Diagram-level form]\label{cor:coreflection-diagrams}
Let $\cat{C}$ be a small category. Post-composition with $\Phi$ and $\Psi$
gives functors $\Phi_{*}, \Psi_{*}$ between the diagram categories
$[\cat{C}, \StatMod]$ and $[\cat{C}, \Statk]$ with $\Phi_{*} \dashv \Psi_{*}$
and $\Psi_{*}\Phi_{*} = \id$. Consequently, taking
$\cat{C} = \catV \times \catD^{\op}$:
\begin{enumerate}
\item Br\o ns' bijection (Proposition~\ref{prop:brons}) transports along
$\Phi_{*}$ to a bijection between McCullagh--Br\o ns models and those functors
$\catV \times \catD^{\op} \to \Statk$ whose two projections are
$D \circ U_\catK\Theta\,\pr_\catO^{\op}\pr_2$ and
$\Gamma(\id \times \pr_\catU^{\op})$; the projection part of
Proposition~\ref{prop:brons} is thus replaced by its discretisation on the
parameter leg.
\item A diagram in $\Statk$ lies in the essential image of $\Phi_{*}$ if and
only if it is pointwise discrete on the parameter side, the obstruction at
each index being the counit component of
Proposition~\ref{prop:coreflection}(4).
\end{enumerate}
In particular, the diagram-level identification holds up to a coreflection, and
pointwise discreteness of the parameter object is not a technical convenience
but exactly its boundary.
\end{corollary}

\begin{proof}
Post-composition with an adjunction is an adjunction between diagram
categories, with unit and counit computed pointwise, and
$\Psi_{*}\Phi_{*} = (\Psi\Phi)_{*} = \id$. (1) A functor
$F \colon \catV \times \catD^{\op} \to \Statk$ with the stated projections has
parameter legs of the form $D\Theta_0$, so $F = \Phi_{*}\Psi_{*}F$ (both sides have the same underlying functions, and a map out of a discrete
space is measurable); hence $F \mapsto \Psi_{*}F$ and $G \mapsto \Phi_{*}G$
are mutually inverse between such functors and the functors
$\catV \times \catD^{\op} \to \StatMod$ with projections
$U_\catK\Theta\,\pr_\catO^{\op}\pr_2$ and $\Gamma(\id \times \pr_\catU^{\op})$,
and Proposition~\ref{prop:brons} identifies the latter with models.
(2) If $F$ is pointwise discrete on the parameter side, each counit component
$(\varepsilon_{\Theta_c},\id)$ is invertible by
Proposition~\ref{prop:coreflection}(4), so $\Phi_{*}\Psi_{*}F \Rightarrow F$
is a natural isomorphism and $F$ lies in the essential image; conversely
$\Phi_{*}G$ is pointwise discrete by construction, and pointwise discreteness
is invariant under isomorphism in $[\cat{C},\Statk]$, since $X \cong DA$ in
$\Meas$ forces $\Sigma_X = 2^{X}$.
\end{proof}

Two consequences for how the open problem should now be stated. First, the
object-level identification of Proposition~\ref{wprop:copower} and the
diagram-level statement are not competing formulations to choose between: they
are the restriction and the transport of one adjunction, and the question of
which is ``correct'' does not arise. Second, the passage from $\StatMod$ to
$\Arrdet{\Stoch}$ factors as $\Jc \cong \Jcm\Phi$, the isomorphism being
the canonical identification $D\Theta_0 \cong \Theta_0 \cdot I$ in $\Stoch$,
and the two factors behave quite differently. The first,
$\Phi \colon \StatMod \to \Statk$, is unconditional. Every hypothesis in
Proposition~\ref{wprop:copower}, the measurable-cardinal condition~(a) and
the countably-generated condition~(b), is used only on the second,
$\Jcm \colon \Statk \to \Arrdet{\Stoch}$. What remains open is
therefore confined to $\Jcm$.

Whether Br\o ns' equivalence (Proposition~\ref{prop:brons}), transported as in
Corollary~\ref{cor:coreflection-diagrams}, continues to hold with
$\Arrdet{\Stoch}$ in place of $\Statk$, and under which hypotheses
$\Jcm$, hence the composite $\Jc \cong \Jcm\Phi$ onto
$\Arrcop{\Stoch}$, is an equivalence, is left open here.

\section{Meaningfulness as naturality}\label{sec:meaningful}

Two criteria are on the table, and it is worth keeping them apart before
merging them. Tjur's discussion~\cite{tjur_discussion_2002}
of~\cite{mccullagh_what_2002} proposes an elementary one: a
parameter function, a rule assigning to each finite sample of covariate
values a parameter of the corresponding model, in his words a ``function of
the unknown distribution''~\cite[p.~1299]{tjur_discussion_2002}, is
\emph{meaningful} when it is invariant under the formation of marginal
models, that is, when reducing the design by removing some of the sampled
$x$'s leaves the associated parameter unchanged. The removal deletes sample
points together with their covariate values, not points of a covariate
space; Lemma~\ref{lem:tjurshadow} identifies its action on the quantities
below with naturality over the insertions. McCullagh's own
criterion~\cite[\S4.5]{mccullagh_what_2002} is that a subparameter be a natural
transformation of functors on the design category.\footnote{\label{fn:indexings}The two indexings
agree on the families compared here, so the comparison below is over a single
class of morphisms. Designs index McCullagh's components in \S4.5 and
natural over all design morphisms $(\psi_{d},\psi_{c})$, whereas
Definition~\ref{def:meaningful} indexes quantities by objects of $\catO$ and
quantifies over its morphisms; on families factoring through $\pr_{\catO}$ the
squares with identity covariate leg are trivial, and the remaining ones are
exactly the morphisms of $\catO$, since for
any $\psi \colon \omega \to \omega'$ and any design $(u,\omega,x)$ the pair
$(\id_{u},\psi) \colon (u,\omega,x) \to (u,\omega',\psi\circ x)$ is a morphism
of $\catD$, so that $\pr_{\catO}$ is surjective on morphisms.} Tjur quantifies over
removal of sample points, acting, by Lemma~\ref{lem:tjurshadow}, through
the insertions alone; McCullagh over all design morphisms. Neither source
asserts that the two coincide as mathematical requirements, and in general they
do not, but how far they
differ depends on a convention, and it is worth being explicit about whose.
McCullagh takes the morphisms of $\catO$ to be
injective~\cite[\S4.1]{mccullagh_what_2002}, so under his own convention
merging is not a morphism at all, and the gap between his criterion and Tjur's
consists exactly of the bijections—that is, of relabelling.
Definition~\ref{def:designdata} follows Br\o ns in imposing no such
restriction, and every map of covariate spaces factors as a surjection
followed by an injection; so in the setting of this paper the three quantifier
classes are nested strictly,
\[
\underbrace{\text{insertions}}_{\text{Tjur}}
\;\subsetneq\;
\underbrace{\text{all injections (insertions and bijections)}}_{\text{McCullagh as stated}}
\;\subsetneq\;
\underbrace{\text{all of } \catO}_{\text{naturality here}},
\]
and the corresponding conditions are ordered the other way, naturality over
all of $\catO$ being the strongest. The outer step is merging; the inner step
is relabelling. Here and throughout, the \emph{insertions} are the proper
inclusions $\omega' \subsetneq \omega$ between objects of $\catO$, the
covariate legs of Tjur's reductions, as fixed in Example~\ref{ex:oneway-setup},
and $\catO_{\mathrm{red}} \subseteq \catO$ denotes the subcategory they
generate, whose morphisms are the inclusions (proper or identity). Insertions
carry the surviving labels unchanged; an injection that is not an inclusion
is a relabelling followed by an insertion, and naturality along it is not
part of Tjur's condition. What the bridge buys is that all three
conditions become
computations in the same ambient category so that the gap can be measured rather
than elided, and, as \S\ref{sec:hochschild} shows for the one-way layout, it
is the reduction morphisms that carry the defect in the motivating example, so
that Tjur's weaker criterion already suffices to detect it there.

Throughout we take \emph{meaningfulness} in the strong (naturality) sense, and
record where the weaker Tjur condition would do. We fix the definition in the
bridged setting.

The value space of a quantity may itself depend on the design, and a change of
design then transports values: in the one-way layout
(Example~\ref{ex:oneway-setup}, $\omega$ a finite label set) the vector of
group means at $\omega'$ is a point of $\R^{\omega'}$, carried along
$\psi \colon \omega \to \omega'$ to $\R^{\omega}$ by the restriction
$\psi^{*}$; in the regression example (Example~\ref{ex:linreg}) the
coefficient vector at $\omega'$ is a point of $\omega'^{*}$, carried to
$\omega^{*}$ by the transpose of $\psi$. The target of a design-indexed
quantity is accordingly a functor $\Lambda \colon \catO^{\op} \to \Meas$, of
the same variance as $\Theta_\Meas$, so that the squares below compose. A
constant $\Lambda$ recovers the familiar fixed-scale case
(Example~\ref{ex:alphabetaxbar} and the whole of \S\ref{sec:hochschild});
Example~\ref{ex:tjurgap} has a genuinely varying target.

\begin{definition}[Design-indexed quantity; meaningfulness]\label{def:meaningful}
Let $\mathbf{P} \colon \mathbf{L} \Rightarrow \mathbf{R}$ be a bridged model
(Theorem~\ref{thm:bridge-meas}). Let
$\Lambda \colon \catO^{\op} \to \Meas$ be a functor, the \emph{target}:
$\Lambda(\omega)$ is the measurable space of values at design $\omega$. A
\emph{design-indexed quantity} for $\mathbf P$ with target
$\Lambda$ is a family of morphisms in $\Stoch$
\[
g_{\omega} \colon \Theta_\Meas(\omega) \rightsquigarrow \Lambda(\omega),
\qquad \omega \in \ob\catO
\]
(the deterministic case is when $g$ is an ordinary parameter functional). The
quantity is \emph{meaningful} if $g$ is a natural transformation
$\del\,\Theta_\Meas \Rightarrow \del\,\Lambda$ (deterministic case) or
$\Theta_\Meas \Rightarrow \Lambda$ read in $\Stoch$ (kernel case): i.e.\ for
every $\psi \colon \omega \to \omega'$ in $\catO$,
\begin{equation}\label{eq:tjursquare}
g_{\omega} \circ \del\Theta_\Meas(\psi) \;=\; \del\Lambda(\psi) \circ
g_{\omega'}
\quad\text{in } \Stoch .
\end{equation}
\end{definition}

A design-indexed quantity $g$ is \emph{Tjur-natural} when~\eqref{eq:tjursquare}
holds for every $\psi$ of $\catO_{\mathrm{red}}$, the subcategory of
inclusions generated by the insertion morphisms as fixed at the opening of
this section; Lemma~\ref{lem:tjurshadow}
warrants the name: the formation of marginal models is removal of sample
points, and on quantities indexed over $\catO$ it acts through the
insertions and through nothing else.

\medskip
\noindent\textbf{Tjur reductions.} Fix design data in $\Set$ as in
Example~\ref{ex:oneway-setup}, with $\catO$ a category of finite sets such as
$\catO_{\mathrm{lab}}$: $\cat{C} = \Set$, both $J$'s the evident
underlying-set functors, and
$\catU$ a full subcategory of $\FinSet_{\mathrm{inj}}$ containing the
underlying sets of the objects of $\catO$ (so that every subsample inclusion
between such sets is a morphism of $\catU$). The regression instantiations
of \S\ref{sec:worked}, whose covariate objects are vector spaces, fall
outside this setting. A \emph{Tjur reduction} of a design
$(u,\omega,x)$ with $x$ surjective is a morphism of $\catD$
\[
(\varphi,\psi) \colon (u',\omega',x') \longrightarrow (u,\omega,x)
\]
in which $\varphi \colon u' \hookrightarrow u$ is a proper subsample
inclusion, $\omega' = x(\varphi(u'))$ the surviving support, $x'$ the
corestriction of $x \circ \varphi$, and
$\psi \colon \omega' \hookrightarrow \omega$ the inclusion: the removal of
the sample points $u \setminus \varphi(u')$ together with the covariate
values no longer observed, Tjur's ``removal of some of the
$x$'s''~\cite[p.~1299]{tjur_discussion_2002}. His condition at such a
reduction compares the parameter assigned to the marginal model with the
parameter for the bigger sample, the former being for him a function of the
marginal distribution. Since a bridged model is natural at $(\varphi,\psi)$,
this is square~\eqref{eq:ND}, the marginal law at the reduced design
depends on $\theta$ only through $\Theta_\Meas(\psi)\theta$; the comparison
is therefore square~\eqref{eq:tjursquare} at $\psi$, and we say $g$
\emph{satisfies Tjur's condition} at the reduction when that square
commutes.

\begin{lemma}[Tjur reductions act through their covariate
shadows]\label{lem:tjurshadow}
In the setting above, a design-indexed quantity $g$ satisfies Tjur's
condition at every Tjur reduction of every design with a surjective design map
if and only if $g$ is Tjur-natural.
\end{lemma}

\begin{proof}
If $g$ is Tjur-natural: the covariate leg $\psi$ of a Tjur reduction is the
inclusion of the surviving support, hence an insertion (a proper inclusion)
when the removal empties some fibre of $x$ and the identity otherwise;
square~\eqref{eq:tjursquare} holds at the former by hypothesis and at the
latter trivially. Conversely, it suffices to check~\eqref{eq:tjursquare} at
each insertion $\psi \colon \omega' \subsetneq \omega$, since such squares
paste along composites and hold trivially at identities. Consider the
tautological design $(\omega,\omega,\id_\omega)$, one unit per treatment
label, with surjective design map, and the subsample
$\varphi = \psi \colon \omega' \hookrightarrow \omega$, a morphism of $\catU$
because $\omega'$ and $\omega$ are objects of $\catO$, hence of $\catU$, and
$\catU$ is full in $\FinSet_{\mathrm{inj}}$. The surviving support is
$x(\varphi(\omega')) = \omega'$, so
$(\psi,\psi) \colon (\omega',\omega',\id) \to (\omega,\omega,\id)$ is a Tjur
reduction with covariate leg $\psi$, and Tjur's condition at it is
square~\eqref{eq:tjursquare} at $\psi$.
\end{proof}

\begin{remark}[What the indexing screens off]\label{rem:tjurscope}
We plainly state two limitations of the Lemma. First, the same
removal recorded without corestriction, as
$(\varphi,\id_\omega) \colon (u',\omega,x\circ\varphi) \to (u,\omega,x)$, has
identity covariate leg, and for quantities indexed over $\catO$ every such
square is trivial (footnote~\ref{fn:indexings}); the force of Tjur's
criterion is carried entirely by the corestriction, that is, by his own
requirement that the parameter of the marginal model be a function of the
marginal distribution, which, for the models of \S\ref{sec:worked}, the
surviving support parametrises. Second, Tjur's parameter functions are
indexed by samples, and before his criterion is imposed they include the
sample size $n$ and the average $\bar x$ weighted by
multiplicities~\cite[p.~1299]{tjur_discussion_2002}; quantities in the sense
of Definition~\ref{def:meaningful} are indexed by $\catO$ and can express
neither, replication-dependence is screened off by the indexing rather
than eliminated by the criterion, and its model-side trace is the
$\catU$-slot failure of Example~\ref{ex:centring}. What the criterion
adjudicates on $\catO$-indexed quantities is dependence on the observed
covariate values, and there, by the Lemma, sample-level invariance and
naturality over $\catO_{\mathrm{red}}$ coincide.
\end{remark}

\begin{proposition}[Meaningfulness is naturality; comparison with Tjur's
criterion; type-level reading]\label{prop:tjur}
Let $g$ be a design-indexed quantity. Then:
\begin{enumerate}
\item if $g$ is deterministic in the sense $g_\omega = \del f_\omega$ for
measurable $f_\omega \colon \Theta_\Meas(\omega) \to \Lambda(\omega)$ (an
ordinary parameter functional), and the $\sigma$-algebras of the
$\Lambda(\omega)$ separate points, then $g$ is meaningful iff the $f_\omega$
form a natural transformation $\Theta_\Meas \Rightarrow \Lambda$ in $\Meas$;
\item[(1$'$)] meaningfulness implies Tjur-naturality; the converse holds when
$\catO_{\mathrm{red}} = \catO$, and fails in general
(Example~\ref{ex:tjurgap});
\item meaningfulness is stable under postcomposition with natural
transformations of targets $\Lambda \Rightarrow \Lambda'$ (in $\Meas$, or in
$\Stoch$ for kernel quantities) and under pullback along functors
$F \colon \catO' \to \catO$ (change of covariate scheme; the bridged model
$\mathbf{P}$ is restricted along $F$ as well);
\item in the functor category
$\Stoch^{\,\catO^{\op}}$ (objects: functors; morphisms: natural
transformations), $g$ is meaningful iff it \emph{is} a morphism
$\del\Theta_\Meas \to \del\Lambda$ (deterministic case) or
$\Theta_\Meas \to \Lambda$ (kernel case). A non-natural family is not a
morphism of the category, it lacks the asserted type.
\end{enumerate}
\end{proposition}

\begin{proof}
(1) With $g = \del f$, square~\eqref{eq:tjursquare} reads
$\del(f_\omega \circ \Theta_\Meas(\psi)) = \del(\Lambda(\psi) \circ
f_{\omega'})$, and $\del$ is faithful on maps into the point-separating
space $\Lambda(\omega)$ (Definition~\ref{def:dirac}), so this holds iff
$f_\omega \circ \Theta_\Meas(\psi) = \Lambda(\psi) \circ f_{\omega'}$.
(1$'$): the implication and the equality case are immediate, Tjur-naturality
being~\eqref{eq:tjursquare} with the quantifier restricted to
$\catO_{\mathrm{red}}$; Example~\ref{ex:tjurgap} supplies a quantity that is
Tjur-natural but not meaningful (Example~\ref{ex:alphabetaxbar}, by contrast,
fails already on $\catO_{\mathrm{red}}$ and does not separate the two
conditions). (2) is
functoriality of pasting. (3) is a restatement of the definition of the
functor category.
\end{proof}

\begin{example}[Tjur's non-example: the design-averaged
expectation]\label{ex:alphabetaxbar}
Recast in the present notation, Tjur's own
non-example~\cite[pp.~1299--1300]{tjur_discussion_2002} is the simple regression
model with $\Theta(\omega) = \R^2 \times \R_{>0}$, coordinates
$(\alpha,\beta,\sigma^2)$, $\catO$ the category of finite covariate
selections, and $\Lambda$ constant at $\R$ (with $\Lambda(\psi) = \id$);
this recasts the unreplicated case, the design map injective, so that
sample points and covariate values coincide, the general case, with
$\bar x$ weighted by multiplicities, is screened off by the
$\catO$-indexing (Remark~\ref{rem:tjurscope}). Put
\[
g_{\omega}(\alpha,\beta,\sigma^2) \;=\; \alpha + \beta\,\bar x_{\omega},
\qquad
\bar x_{\omega} = |\omega|^{-1}\textstyle\sum_{x \in \omega} x .
\]
Each $g_\omega$ is a legitimate parameter of the model at $\omega$. But for
an insertion $\psi \colon \omega \subsetneq \omega'$ the parameter map
$\Theta(\psi)$ is the identity, while $\bar x_\omega \neq \bar x_{\omega'}$ in
general, so~\eqref{eq:tjursquare} reads $\alpha + \beta\bar x_{\omega} = \alpha
+ \beta\bar x_{\omega'}$ and fails off the locus $\beta = 0$. The failure is
already visible on $\catO_{\mathrm{red}}$, so the quantity is non-meaningful in
both Tjur's sense and ours. The same argument, with $\Theta(\psi)$ again the
identity, disqualifies the correlation coefficient, whose design dependence
enters through the covariate
variance~\cite[\S6.2]{mccullagh_what_2002}; here it is
the design-\emph{scale} rather than the design-\emph{mean} that fails to
transport. In the bridged setting both failures are equalities that fail,
not methodological unease.
\end{example}

\begin{example}[A quantity natural on $\catO_{\mathrm{red}}$ but not on
$\catO$]\label{ex:tjurgap}
Let $\catO = \catO_{\mathrm{lab}}$ (Example~\ref{ex:oneway-setup}), let
$\Theta$ be any parameter functor, and take as target
$\Lambda(\omega) = \R^{\omega \times \omega}$ with
$\Lambda(\psi) = (\psi\times\psi)^{*}$, i.e.\
$\Lambda(\psi)(v)_{(a,b)} = v_{(\psi a,\,\psi b)}$. Put
$g_\omega(\theta)_{(a,b)} = [a \neq b]$, the indicator that two treatment
labels are distinct. Square~\eqref{eq:tjursquare} at
$\psi \colon \omega \to \omega'$ reads $[a \neq b] = [\psi a \neq \psi b]$
for all $a,b \in \omega$, which holds iff $\psi$ is injective. Hence $g$ is
natural along every injection, in particular along every insertion, hence
Tjur-natural, and natural even under McCullagh's criterion as stated (all
injections, i.e.\ insertions and bijections), and fails along every merge, for instance
along $p \colon \{1,2\} \to \{*\}$: the gap between Tjur-naturality and
meaningfulness is realised at the merges.
\end{example}

\begin{corollary}[Finite verification]\label{cor:decidable}
Suppose $\catO$ has finitely many morphisms, $\Theta_\Meas(\omega)$ and
$\Lambda(\omega)$ are finite (discrete) for every $\omega$, and the entries of
the kernels $g_\omega$ lie in a subfield $k \subseteq \R$ with decidable
equality (for instance $\mathbb{Q}$, or the real algebraic numbers). Then
meaningfulness of $g$ is decidable: it is equivalent to the finite
conjunction, over the morphisms $\psi \colon \omega \to \omega'$ of any
generating set of $\catO$, of the matrix equalities~\eqref{eq:tjursquare},
each between two $|\Theta_\Meas(\omega')| \times |\Lambda(\omega)|$ stochastic
matrices with entries in $k$.
\end{corollary}

\begin{proof}
If~\eqref{eq:tjursquare} holds at $\psi$ and at $\psi'$ then it holds at
$\psi'\psi$, by pasting the two squares ($\del\Theta_\Meas$ and $\del\Lambda$
being functors), so it suffices to check a generating set, which is finite;
each check is an equality of finitely many elements of $k$, decidable by
hypothesis.
\end{proof}

\section{A cohomological formulation of coherence for finite
designs}\label{sec:hochschild}

This section formulates, it does not yet fully prove, the connection
between meaningfulness and the cohomology of diagrams of algebras, which
links the present framework to the diagrammatic
Hochschild/\allowbreak Gerstenhaber--Schack theory
of~\cite{gerstenhaber_schack_1983,baues_cohomology_1985}, as revisited
in~\cite[\S3]{caputi_diagrammatic_2025}. The logical
role of this section is precisely
delimited: cohomology here \emph{controls obstructions} to coherence; it is
upstream of the bridge, not a proof of it. A distinct cohomological tradition
attaches invariants to statistical structures directly; its relation to the
complex constructed here is set out in Remark~\ref{rem:infocohomology} at the
end of the section.

\subsection{The diagram of algebras of a finite design scheme}

Let $\catO$ be a finite category of covariate schemes and let all spaces,
the parameter sets $\Theta(\omega)$ and the target sets $\Lambda(\omega)$,
be finite, as in Corollary~\ref{cor:decidable}. Fix a field $k \supseteq \mathbb{Q}$.

\begin{definition}\label{def:diagalg}
Write $\Alg_k$ for the category of associative unital $k$-algebras and unital
$k$-algebra homomorphisms. The \emph{parameter diagram of algebras} is the
functor
\[
A \colon \catO \longrightarrow \Alg_k,
\qquad
A(\omega) \coloneqq k^{\,\Theta(\omega)}
\]
(functions on the parameter set, a commutative $k$-algebra under pointwise
operations), covariant because $\Theta$ is contravariant and
$k^{(-)}$ is contravariant again: $A(\psi) = (\Theta(\psi))^{*}$ is
precomposition. Similarly, each target functor $\Lambda$ yields
$B(\omega) = k^{\Lambda(\omega)}$.
\end{definition}

A \emph{deterministic} design-indexed quantity $g = \{g_\omega \colon
\Theta(\omega) \to \Lambda(\omega)\}$ dualises to a family of algebra maps
$g^{*} = \{g_\omega^{*} \colon B(\omega) \to A(\omega)\}$, and $g$ is
meaningful (Definition~\ref{def:meaningful}) iff $g^{*}$ is a natural
transformation of diagrams $B \Rightarrow A$: all spaces being finite and
discrete, meaningfulness is naturality of the underlying maps
(Proposition~\ref{prop:tjur}(1)), and $k^{(-)}$ is faithful, a map $f$ is
recovered from $f^{*}$ through the indicator functions of its fibres, so
naturality of $g$ and of $g^{*}$ are equivalent. Failure of naturality is
measured, for each $\psi \colon \omega \to \omega'$, by the difference
\begin{equation}\label{eq:defect}
c(\psi) \;\coloneqq\; A(\psi)\, g_{\omega}^{*} \;-\; g_{\omega'}^{*}\, B(\psi)
\;\in\; \Hom_k\bigl(B(\omega), A(\omega')\bigr),
\end{equation}
a $1$-cochain on $\catO$ in the sense of
Baues--Wirsching~\cite[(1.4)]{baues_cohomology_1985}, with coefficients in
the natural system $D \colon F\catO \to \mathrm{Ab}$ obtained by pulling
back the bimodule $M(\omega,\omega') = \Hom_k(B(\omega),A(\omega'))$
along the forgetful functor
$F\catO \to \catO^{\op} \times \catO$~\cite[(1.16)--(1.18)]{baues_cohomology_1985}:
$D(\psi) = \Hom_k(B(\omega),A(\omega'))$ for $\psi \colon \omega \to
\omega'$, a morphism $(\alpha,\beta) \colon \psi \to \alpha\psi\beta$ of
the factorization category acting by $f \mapsto A(\alpha)\,f\,B(\beta)$. With
bimodule coefficients this is Hochschild--Mitchell
cohomology~\cite[(8.5)]{baues_cohomology_1985}; when
$B = A$ this complex is the Hochschild-degree-$1$ row of the Gerstenhaber--Schack
bicomplex of the diagram $A$~\cite[Prop.~3.8]{caputi_diagrammatic_2025} (defined
over a poset in~\cite{gerstenhaber_schack_1983}, with diagrams over a small
category as in~\cite{gerstenhaber_simplicial_1983}). The group $H^{1}$
computed below is the cohomology of that row for its horizontal,
Baues--Wirsching, differential, i.e.\ it lives in bidegree $(1,1)$ of the
bicomplex; it is not Gerstenhaber--Schack cohomology in total degree~$2$, and
no vertical (Hochschild) differential is used here. The relation of the
bicomplex to Baues--Wirsching cohomology is controlled by a spectral sequence
attributed to Robinson~\cite{robinson_cohomology_2008}
in~\cite[\S1, Thm.~3.11]{caputi_diagrammatic_2025}.

\begin{remark}[Why $H^{1}$ is the right home]\label{rem:endH0}
The cohomological formulation is forced rather than chosen. Meaningfulness of
a design-indexed quantity is the assertion that the family $g^{*}$ is an
element of $\Nat(B,A) = \int_{\omega} \Hom_k\bigl(B(\omega),A(\omega)\bigr)$,
an end (the subgroup of $\prod_\omega \Hom_k(B(\omega),A(\omega))$ cut out
by the equations $A(\psi)h_\omega = h_{\omega'}B(\psi)$, one for each
$\psi$;~\cite[IX.5]{maclane_categories_1998}); and $H^{0}$ of the
Baues--Wirsching complex of the diagram \emph{is} that end. The defect
cochain~\eqref{eq:defect} is then a $1$-cocycle (Lemma~\ref{lem:cocycle})
whose class in the first derived functor of that
end~\cite[(8.5)]{baues_cohomology_1985}, $H^{1}$, is the obstruction to correcting $g^{*}$ by an arbitrary $0$-cochain,
so $H^{1}$ is the only place where a non-meaningful family could leave a trace,
and the class-relativity of
Proposition~\ref{wprop:obstruction}(2) below is the statement that the
$0$-cochain group is too large for that trace to survive unless the admissible
correction class is cut down.
\end{remark}

\begin{lemma}[Cocycle property]\label{lem:cocycle}
The cochain $c$ of~\eqref{eq:defect} satisfies the $1$-cocycle identity
$c(\psi'\psi) = A(\psi')\,c(\psi) + c(\psi')\,B(\psi)$ for composable
$\psi,\psi'$, and $c = 0$ iff $g$ is meaningful.
\end{lemma}

\begin{proof}
Expand $c(\psi'\psi)$ using functoriality of $A, B$ and add-subtract the
mixed term $A(\psi')g^{*}_{\omega'}B(\psi)$:
\[
A(\psi'\psi)g^{*}_{\omega} - g^{*}_{\omega''}B(\psi'\psi)
= A(\psi')\bigl(A(\psi)g^{*}_{\omega} - g^{*}_{\omega'}B(\psi)\bigr)
+ \bigl(A(\psi')g^{*}_{\omega'} - g^{*}_{\omega''}B(\psi')\bigr)B(\psi).
\]
The second statement is the equivalence recorded after
Definition~\ref{def:diagalg}: $c = 0$ says that $g^{*}$ is natural, which
holds iff $g$ is.
\end{proof}

\begin{proposition}[Obstruction formulation]\label{wprop:obstruction}
Let $g$ be a design-indexed quantity given only on the objects of $\catO$
(a ``raw family''), fix an \emph{admissible class} $\mathcal{C} \subseteq
C^{0} = \prod_\omega \Hom_k(B(\omega),A(\omega))$ of $0$-cochains, and
suppose one seeks a correction $h \in \mathcal{C}$ such
that the corrected family $g^{*} + h$ is natural. (A $0$-cochain is merely a
family of linear maps; $g^{*}+h$ is the dual of a design-indexed quantity only
when it is a family of algebra homomorphisms, which $\mathcal{C}$ must ensure,
the recalibration class of Example~\ref{ex:oneway} does so by
construction.) Then:
\begin{enumerate}
\item a corrected natural family exists iff the equation $c = -\,d^{0}h$ is
solvable with $h \in \mathcal{C}$; for a linear subspace $\mathcal{C}$ this
is finite linear algebra, and for a finite non-linear class (such as the
recalibrations of Example~\ref{ex:oneway}) it is finite combinatorics. Both $c$
and the correcting $h$, when it exists, are computed exactly;
\item the obstruction is intrinsically \emph{class-relative}. When corrections
range over the whole $0$-cochain group, the obstruction vanishes identically and
carries no information: by~\eqref{eq:defect} the defect cochain of any raw
family is by construction the coboundary $d^{0}g^{*}$ of that family, so
$[c_g] = 0$ in $H^{1}$ for every $g$ whatever, with no hypothesis. Admissible
classes must therefore be \emph{proper} subsets of the $0$-cochains to carry
any obstruction at all;
\item $H^{1}(\catO;\Hom(B,A))$ is accordingly not an obstruction group for
individual quantities but an invariant of the design scheme (and of the
target $\Lambda$), computable by finite linear algebra;
Example~\ref{ex:oneway-h1} records a nonzero instance.
\end{enumerate}
\end{proposition}

\begin{remark}[Why part (2) matters]
The identity $c_g = d^{0}g^{*}$ is the reason the naive reading of part~(1),
``non-meaningfulness is certified by a nonzero class in $H^{1}$'', is
false as stated, and the phrase ``complete obstruction in $H^{1}$'' should not
be used without naming the class. Equivalently: with corrections ranging over
all $0$-cochains, the zero family is natural and corrects anything. The cocycle
identity of Lemma~\ref{lem:cocycle} is then the identity $d^{1}d^{0} = 0$.
\end{remark}

\begin{proof}
If $g^{*}+h$ is natural for some $0$-cochain $h$ then $c = -\,d h$, where $d$
is the Baues--Wirsching differential $\,(d h)(\psi) = A(\psi) h_\omega -
h_{\omega'} B(\psi)$, so $[c]=0$ in $H^{1}$, a statement carrying no
information about $\mathcal{C}$; conversely if $c = -\,d h$ with
$h \in \mathcal{C}$ then $g^* + h$ is natural by reversing the computation.
Part~(2) is the
substitution $h = -g^{*}$ in this equivalence, and part~(3) follows since the
complex is finite-dimensional in each degree.
\end{proof}

The linear-algebraic skeleton (Lemma~\ref{lem:cocycle} and the displayed
equivalence) is elementary, and the proof above establishes all three parts.
The statement is discharged, in the class-relative form that part~(2) forces, by
an exact computation for the one-way layout; the findings
are recorded in Example~\ref{ex:oneway} and its successors.

\subsection{The one-way layout as test case}

\begin{example}[One-way layout: the full finite complex]\label{ex:oneway}
Truncate Example~\ref{ex:oneway-setup} to the smallest nontrivial covariate
category. \emph{The choice of that category is not a matter of convenience:
the merge morphism alone cannot host the test case.} Write
$\catO_{\mathrm{arr}}$ for the two-object category with $\omega_2 = \{1,2\}$
(two treatment groups), $\omega_1 = \{*\}$ (groups merged) and the single
non-identity morphism $p \colon \omega_2 \to \omega_1$, and
$\catO_{\mathrm{full}}$ for the full subcategory of $\FinSet$ on the same two
objects, which adds the injections $i_1, i_2 \colon \omega_1 \to \omega_2$
(subgroup selection), the swap $\sigma$ (treatment relabelling) and the
constants $c_k = i_k p$: eight morphisms in all, with $p\,i_k = \id$ and
$\sigma i_1 = i_2$. Identifying $\omega_1 = \{*\}$ with $\{1\}$, $i_1$ is
the insertion $\{1\} \subsetneq \{1,2\}$, the shadow of Tjur's reduction
to the first group, Lemma~\ref{lem:tjurshadow}, and $i_2 = \sigma i_1$
is that insertion followed by a relabelling; the reduction to the second
group is the insertion $\{2\} \subsetneq \{1,2\}$, which $i_2$ represents up
to the same identification. Discretise the parameters:
$\Theta(\omega) = M^{\omega} \times S$ with $M, S$ finite sets of admissible
means and dispersions, $\Theta(\psi)(\mu, s) = ((\mu,\mu), s)$ the diagonal.
Then:
\[
A(\omega_1) = k^{M \times S}, \qquad
A(\omega_2) = k^{M^{2} \times S}, \qquad
A(\psi) = \text{restriction to the diagonal}.
\]
The Baues--Wirsching $1$-cochains with coefficients as above are indexed by
the single non-identity morphism $\psi = p$ (identities carry no information
after normalisation), so
\[
C^{0} = \Hom_k(B(\omega_1), A(\omega_1)) \oplus \Hom_k(B(\omega_2),
A(\omega_2)),
\qquad
C^{1}_{\mathrm{norm}} = \Hom_k\bigl(B(\omega_2), A(\omega_1)\bigr),
\]
the latter being the coefficient of the natural system at $\psi = p$, and
$H^{1}$ is the cokernel of the explicit linear map
$d^{0}(h_1,h_2) = A(\psi)h_2 - h_1 B(\psi)$. Everything is therefore a finite
matrix computation once the candidate $g$ is fixed. Take for $g$ the
\emph{marginal} (total) dispersion under balanced allocation,
\[
g_{\omega_1}(\mu,s) = s,
\qquad
g_{\omega_2}\bigl((\mu_1,\mu_2),s\bigr) = s + \tfrac{(\mu_1-\mu_2)^2}{4},
\]
the one-way analogue of Tjur's $\alpha + \beta\bar x$ of
Example~\ref{ex:alphabetaxbar}: a genuine parameter contaminated by a term that
depends on which groups the design includes. The quantity is ours: it should
not be confused with the overdispersion passage of Tjur's discussion, which
concerns whether a generalised linear model admits a scale extension leaving
the link-linear estimates unchanged, a non-existence claim about a family
of models, not a failure of naturality, and one on which nothing here bears.
Throughout, the target
functor $\Lambda$ is constant at the finite set
$T \coloneqq g_{\omega_2}\bigl(\Theta(\omega_2)\bigr) \subset \mathbb{Q}$ of
values of $g$ (which contains $S = g_{\omega_1}(\Theta(\omega_1))$), with
$\Lambda(\psi) = \id_{T}$; accordingly $B(\omega) = k^{T}$ for both objects and
$B(\psi) = \id$. The computation was
carried out in exact arithmetic, the defect cochain, the cocycle identity,
and the corrected linear system over $\mathbb{Q}$; the two cochain-complex
ranks modulo three $31$-bit primes, in agreement, then re-derived exactly over
$\mathbb{Q}$, and returns three findings.
\end{example}

\begin{example}[Findings]\label{ex:oneway-findings}
\emph{(i) The arrow category is blind.} Over $\catO_{\mathrm{arr}}$ the family
above is \emph{natural}: its defect cochain vanishes identically, because the
between-group term $(\mu_1-\mu_2)^2/4$ vanishes on the diagonal sub-model
$\Theta(p)$. Moreover $H^{1}(\catO_{\mathrm{arr}};\Hom(B,A)) = 0$ (for $|M| =
|S| = 2$: $\dim C^{1} = \dim Z^{1} = \dim B^{1} = 16$), so no quantity whatever
carries a class there. It is the subgroup-selection morphisms $i_1, i_2$, not
the merge, that detect the failure, which is to say that the failure is
already visible to Tjur's own criterion, quantified over removal of sample
points alone: the defect at the insertion $i_1$ is a violation of his
condition at the reduction to the first group (Lemma~\ref{lem:tjurshadow}).

\emph{(ii) Non-meaningfulness certified; a deterministic repair recovered.} Over
$\catO_{\mathrm{full}}$ with $M = \{0,1\}$, $S = \{1,2\}$, so $T =
\{1,\tfrac54,2,\tfrac94\}$: the defect cochain is nonzero exactly on
$\{i_1,i_2,c_1,c_2\}$ and zero on $p$ and $\sigma$, and the cocycle identity of
Lemma~\ref{lem:cocycle} was checked entrywise on all $22$ composable pairs.
Solving $c = -\,d^{0}h$ over the full $0$-cochain group, which by
Proposition~\ref{wprop:obstruction}(2) must succeed, returns a corrected
family that is again deterministic; the solution obtained by setting the free
parameters of the row-reduced system to zero is the \emph{within-group}
dispersion $((\mu_1,\mu_2),s) \mapsto s$, i.e.\ the between-group term is
discarded. The solution set is the affine space  $\,-g^{*} + \Nat(B,A)$ and
contains other deterministic members, so no canonicity is claimed for this
one; the vanishing of the unrestricted class (part~(2) of the Proposition)
is thus not a defect of the theory, and the statistically sensible repair is
among the solutions.

\emph{(iii) Nonzero obstruction relative to an admissible class.} Take as
admissible the \emph{information-bearing recalibrations}: corrected quantities
$r_\omega \circ g_\omega$ with $r_\omega \colon T \to T$ and $r_{\omega_2}$
non-constant, so that the quantity may be re-expressed on its value scale but
must remain a non-trivial function of the observed marginal dispersion.
Naturality along $i_1, i_2$ reads $r_{\omega_1}(s) = r_{\omega_2}\bigl(s +
(\mu_1-\mu_2)^2/4\bigr)$ for all $\mu, s$, and forces $r_{\omega_2}$ to be
constant on the partition of $T$ generated by $s \sim s + (\mu_1-\mu_2)^2/4$,
reducing feasibility to an exact union--find computation; the conditions
along the other morphisms add nothing (naturality along $\sigma$ holds for
every recalibration since $g_{\omega_2}$ is symmetric; along $p$ it reads
$r_{\omega_2} = r_{\omega_1}$ on $S$, the case $\mu_1 = \mu_2$ of the
condition along $i_k$; and along $c_k = i_k p$ it follows from the latter).
Call the design \emph{scale-confounded} when some marginal value
$s + (\mu_1-\mu_2)^2/4$ with $\mu_1 \neq \mu_2$ is itself an admissible
within-group dispersion, i.e.\ lies in $S$. With $M = \{0,1\}$ and
$S = \{1,\tfrac54\}$, or already with $S = \{1,2\}$ at three mean levels
$M = \{0,1,2\}$, the partition collapses to a single block, every natural
recalibration is constant, and the obstruction relative to this class is
nonzero: no information-bearing recalibration of the marginal dispersion is
natural. (With $M = \{0,1\}$, $S = \{1,2\}$ as in~(ii) the partition has the
two blocks $\{1,\tfrac54\}$ and $\{2,\tfrac94\}$, and a non-constant natural
recalibration exists.)
\end{example}

\begin{example}[A new invariant]\label{ex:oneway-h1}
Computing $H^{1}(\catO;\Hom(B,A))$ for its own sake, part~(3) of
Proposition~\ref{wprop:obstruction}, gives $0$ for $\catO_{\mathrm{arr}}$,
and for $\catO_{\mathrm{full}}$ a group that vanishes at two levels of the
treatment mean but not at three: with $M = \{0,1,2\}$, $S = \{1,2\}$ and
$\Lambda$ constant at $T$ as above (so $|T| = 5$),
$\dim_{k} H^{1}(\catO_{\mathrm{full}};\Hom(B,A)) = 10 = 2\,|T|$, computed
as in Example~\ref{ex:oneway}. The group depends on $\Lambda$ as well as on
$\catO$: with $B$ constant and $B(\psi) = \id$ it is the $|T|$-fold direct sum
of the cohomology with coefficients $\psi \mapsto A(\operatorname{cod}\psi)$.
This is the instance behind the phrase ``nontrivial'' in that part. Its
statistical meaning is not known to us and is not pursued here.
\end{example}

\begin{caveat}[Scope of the cohomological claim]
The correct summary of this section is: \emph{for finite design schemes,
meaningfulness of a design-indexed quantity is decidable, and its correctability
within a prescribed statistically admissible class is controlled by a computable
obstruction relative to that class}. The class-relativity is not a hedge but a
theorem (Proposition~\ref{wprop:obstruction}(2)). This must not be conflated
with a cohomological proof of the bridge: the bridge
(Section~\ref{sec:bridge}) is proved by Kleisli transposition and
needs no cohomology; the cohomology enters one level up, where coherence of
quantities over designs is at stake.
\end{caveat}

\begin{remark}[Information cohomology]\label{rem:infocohomology}
A separate cohomological tradition attaches invariants to statistical
structures directly: entropy is characterised as the essentially unique
$1$-cocycle of information cohomology in the finite discrete
case~\cite{baudot_homological_2015} (for the continuous extension
see~\cite{vigneaux_information_2021}), as the unique (up to a non-negative
scalar) functorial, convex-linear and continuous information-loss
measure~\cite{baez_characterization_2011}, as an operad
derivation~\cite{bradley_entropy_2021}, and defined, in divergence-enriched
Markov categories, as the divergence of a morphism from
determinism~\cite{perrone_markov_2024}; for the functional equation
underlying the cohomological characterisation see~\cite{bennequin_functional_2020}, and
for cohomological models of information networks in the same tradition
see~\cite{manin_homotopy_2024}. The complex of
Definition~\ref{def:diagalg} and the information-cohomological complexes have
different coefficients and different indexing categories, and \emph{no
comparison map between them is constructed here}.

A prior question sits underneath any such comparison, and we state it
separately because it is not cohomological. Information cohomology carries
its own account of conditioning, organised around Shannon's chain
rule~\cite{baudot_homological_2015,vigneaux_information_2021}, and touches
statistical independence only through the vanishing of mutual
information~\cite{baudot_homological_2015}; that account is not the one native
to Markov categories, of which Section~\ref{sec:markov} recalls only the
fragment used here: there, conditioning is disintegration in the sense
of~\cite{cho_disintegration_2019} (Bayesian inversion being the special case
with a prior) and independence is the diagrammatic
condition of~\cite{fritz_synthetic_2020}. Whether the two agree, and under
what hypotheses, is to our knowledge not known; no translation between them
is constructed or assumed anywhere in this paper, and none of its results
depends on one existing. We record this because the comparison of the two
formalisms at the level of these primitives is logically prior to any
comparison of their cohomologies.
\end{remark}


\section{Scope and non-claims}\label{sec:nonclaims}

For the record, two statements adjacent to this paper's results are
formulable in the present framework but are \emph{not} established here or,
to our knowledge, anywhere.

\begin{caveat}[Paradigms as morphism-shapes]
Fisherian, Bayesian, nonparametric, statistical-learning and hybrid
latent-variable modelling can each be \emph{presented} as a choice of
morphism-shape on the shared parameter--data structure of
Section~\ref{sec:bridge} (e.g.\ a latent-variable model as a composite
through an unobserved object). These presentations organise the paradigms;
they do not prove any equivalence between paradigm-specific theorems, and no
such claim is made.
\end{caveat}


\begin{caveat}[Nonparametric and infinite-dimensional limits]
All theorems above are stated for the parametric, measurably-regular case.
Presenting nonparametric models as $(\text{co})$limits of finite-dimensional
ones in $\Stoch$ is a programme, not a technique available off the shelf:
$\Stoch$ is not complete or cocomplete in the ways naive arguments require,
and nothing here should be described as securing the nonparametric case.
\end{caveat}




\appendix
\section{Proof of Proposition~\ref{wprop:copower}}\label{app:copower}

Throughout, for a set $A$ we write $A \cdot I$ for the measurable space
$(A, 2^{A})$. This is the coproduct in $\Meas$ of $A$ copies of the one-point
space, and $\del$ preserves coproducts (Remark~\ref{rem:repr}), so it is the
copower of the monoidal unit in $\Stoch$; with this convention,
``the domain of $q$ is a copower'' means that the domain of $q$ carries its
full power set. For $p \colon \Theta_0 \to \Prb(S)$ the transpose
$\widehat{p} \colon \Theta_0 \cdot I \rightsquigarrow S$ is the kernel
$\widehat{p}(\theta)(B) = p(\theta)(B)$, measurable in $\theta$ because the
domain is discrete; conversely every kernel $q \colon A \cdot I
\rightsquigarrow S$ is $\widehat{\check q}$ for the function
$\check q \colon a \mapsto q(a)$. This is~\eqref{eq:copower}, and the two
assignments are mutually inverse.

The proof uses four facts about deterministic morphisms and one about
$\{0,1\}$-valued measures on power sets. The first is recorded in
\S\ref{sec:markov} with a reference; the short proofs are included to keep
the appendix self-contained.

\begin{lemma}\label{lem:app-det}
Let $k \colon X \rightsquigarrow Y$ be a morphism of $\Stoch$.
\begin{enumerate}
\item $k$ is deterministic if and only if $k(x)(B) \in \{0,1\}$ for every
$x \in X$ and every $B \in \Sigma_Y$~\cite[Ex.~10.4]{fritz_synthetic_2020}.
\item If $\Sigma_Y$ separates points, then $\del$ is injective on
$\Meas(X,Y)$.
\item If $\Sigma_Y$ is countably generated and separates points, then every
$\{0,1\}$-valued probability measure on $Y$ is $\delta_y$ for a unique
$y \in Y$; consequently every deterministic $k \colon X \rightsquigarrow Y$
is $\del f$ for a unique measurable $f \colon X \to Y$.
\item If $v \colon A \rightsquigarrow X$ and $u \colon X \rightsquigarrow A$
are mutually inverse and $v$ is deterministic, then so is
$u$~\cite[Lem.~10.9(c)]{fritz_synthetic_2020}.
\end{enumerate}
\end{lemma}

\begin{proof}
(1) In $\Stoch$ the copy map is the Dirac kernel of the diagonal, so
by~\eqref{eq:push} and~\eqref{eq:pull} the kernel $\mathrm{copy}_Y \circ k$
sends $x$ to the image of $k(x)$ under the diagonal $Y \to Y \times Y$, while
$(k \otimes k) \circ \mathrm{copy}_X$ sends $x$ to the product measure
$k(x) \otimes k(x)$. On a measurable rectangle $B \times C$ these two
measures take the values $k(x)(B \cap C)$ and $k(x)(B)\,k(x)(C)$
respectively, and since rectangles form a $\pi$-system generating
$\Sigma_Y \otimes \Sigma_Y$, two probability measures on $Y \times Y$ agreeing
on rectangles agree. Hence $k$ is deterministic iff
$k(x)(B \cap C) = k(x)(B)\,k(x)(C)$ for all $x$, $B$, $C$. Taking $B = C$
shows that each $k(x)$ is $\{0,1\}$-valued. Conversely, if $\mu$ is a
$\{0,1\}$-valued probability measure then $\mu(B \cap C) = 1$ iff
$\mu(B) = \mu(C) = 1$, if both are $1$ then
$\mu\bigl(Y \setminus (B \cap C)\bigr) \le \mu(Y \setminus B) + \mu(Y
\setminus C) = 0$, and if one is $0$ then $\mu(B \cap C) = 0$, that is,
$\mu(B \cap C) = \mu(B)\,\mu(C)$.

(2) $\del f = \del g$ means $\delta_{f(x)} = \delta_{g(x)}$ for every $x$,
and if $\Sigma_Y$ separates points then $\delta_y = \delta_{y'}$ forces
$y = y'$.

(3) Let $(B_n)_{n}$ generate $\Sigma_Y$ and let $\mu$ be $\{0,1\}$-valued.
Put $C_n = B_n$ if $\mu(B_n) = 1$ and $C_n = Y \setminus B_n$ otherwise, and
$C = \bigcap_n C_n$; then $\mu(C) = 1$, so $C \neq \emptyset$. If
$y, y' \in C$, the class of those $B \in \Sigma_Y$ containing both or neither
of $y, y'$ is a $\sigma$-algebra containing every $B_n$, hence all of
$\Sigma_Y$; as $\Sigma_Y$ separates points, $y = y'$. Thus $C = \{y\}$, so
$\mu(\{y\}) = 1$ and $\mu = \delta_y$, with $y$ unique because $\Sigma_Y$
separates points. For the consequence: by~(1) each $k(x)$ is
$\{0,1\}$-valued, hence $\delta_{f(x)}$ for a unique $f(x)$; $f$ is
measurable because $f^{-1}(B) = \{x : k(x)(B) = 1\} \in \Sigma_X$ for every
$B \in \Sigma_Y$; and $k = \del f$ by definition of $\del$. Uniqueness of $f$
is~(2).

(4) This holds in any Markov category. From $v \circ u = \id_X$ and
determinism of $v$,
$\mathrm{copy}_X = \mathrm{copy}_X \circ v \circ u = (v \otimes v) \circ
\mathrm{copy}_A \circ u$, hence
$(u \otimes u) \circ \mathrm{copy}_X = \bigl((u \circ v) \otimes (u \circ
v)\bigr) \circ \mathrm{copy}_A \circ u = \mathrm{copy}_A \circ u$ by
$u \circ v = \id_A$.
\end{proof}

\begin{lemma}[Ulam]\label{lem:app-ulam}
Call a countably additive $\{0,1\}$-valued probability measure on
$(A, 2^{A})$ that is not a Dirac measure a \emph{two-valued measure} on the
set $A$.
\begin{enumerate}
\item Two-valued measures on $A$ correspond bijectively to countably
complete non-principal ultrafilters on $A$~\cite[(10.3)--(10.4)]{jech_set_2003}.
\item If $A$ carries no two-valued measure, then neither does any set of
cardinality at most $|A|$, nor the power set
$2^{A}$~\cite[Bem.~2 and Satz~1]{ulam_masstheorie_1930}. In particular no
set of cardinality at most that of the continuum carries one.
\item In modern terms, $A$ carries a two-valued measure if and only if there
is a measurable cardinal $\kappa \le |A|$~\cite[Lemma~10.2 and
Def.~10.3]{jech_set_2003}; Ulam's original form is that $A$ carries none
when no weakly inaccessible cardinal is
$\le |A|$~\cite[Satz~(A)]{ulam_masstheorie_1930}, cf.~\cite[Thm.~10.1]{jech_set_2003}.
\end{enumerate}
\end{lemma}

\begin{proof}
(1) For a $\{0,1\}$-valued $\mu$ the class
$\mathcal U = \{B \subseteq A : \mu(B) = 1\}$ is an ultrafilter on $A$,
countably complete by countable additivity, and non-principal exactly when
$\mu$ is not Dirac; conversely a countably complete non-principal ultrafilter
$\mathcal U$ defines such a $\mu$ by $\mu(B) = 1$ iff $B \in \mathcal U$.

(2) If $i \colon B \to A$ is injective and $\mu$ is a two-valued measure on
$B$, then $C \mapsto \mu(i^{-1}(C))$ is a two-valued measure on $A$: it is
countably additive, gives $A$ mass $\mu(B) = 1$, and vanishes on singletons
because $i^{-1}(\{a\})$ has at most one element. This proves the first
assertion. For the power set, identify $2^{A}$ with the set of functions
$A \to \{0,1\}$ and suppose $\nu$ is a two-valued measure on it. For each
$a \in A$ exactly one of the two sets $\{g : g(a) = 0\}$, $\{g : g(a) = 1\}$
has $\nu$-measure $1$; call it $E_a$. We claim that
$\nu\bigl(\bigcap_{a} E_a\bigr) = 1$. Otherwise $\nu\bigl(\bigcup_a (2^{A}
\setminus E_a)\bigr) = 1$; well-order $A$ as $(a_\xi)_{\xi < \lambda}$ and
put $D_\xi \coloneqq (2^{A} \setminus E_{a_\xi}) \setminus
\bigcup_{\eta < \xi} (2^{A} \setminus E_{a_\eta})$, pairwise disjoint sets of
$\nu$-measure $0$ with union $\bigcup_a (2^{A} \setminus E_a)$. Then
$C \mapsto \nu\bigl(\bigcup_{\xi \in C} D_\xi\bigr)$ is a two-valued measure
on $\lambda$, hence on $A$ by the first assertion, contrary to hypothesis.
But $\bigcap_a E_a$ is the single function $g_0$ with $g_0(a)$ the value
singled out by $E_a$, so $\nu(\{g_0\}) = 1$, contradicting non-principality.
Finally, a countable set carries no two-valued measure, since countable
additivity and vanishing on singletons would force total mass $0$; so
neither does $2^{\aleph_0}$, nor any set of cardinality at most
$2^{\aleph_0}$.

(3) If $A$ carries a countably complete non-principal ultrafilter, so does
the cardinal $|A|$, by transport along a bijection; the least cardinal
$\kappa$ carrying one is then $\le |A|$, and it is measurable because such
an ultrafilter on it is automatically
$\kappa$-complete~\cite[Lemma~10.2]{jech_set_2003}. Conversely a measurable
$\kappa \le |A|$ carries a $\kappa$-complete, hence countably complete,
non-principal ultrafilter, which transports to $A$ by the injection argument
of~(2). Ulam's form is subsumed by the modern one, since every measurable
cardinal is strongly inaccessible~\cite[Lemma~10.4]{jech_set_2003}.
\end{proof}

Hypothesis~(a) of Proposition~\ref{wprop:copower} is thus the statement that
no measurable cardinal is $\le |\Theta_0'|$. By Lemma~\ref{lem:app-ulam}(2)
it holds outright, with no large-cardinal input, whenever $\Theta_0'$ has at
most the cardinality of the continuum, which covers every parameter set in
this paper, and it holds for every set if measurable cardinals do not
exist.

\begin{proof}[Proof of Proposition~\ref{wprop:copower}]
\emph{The functor.} On objects, $\Jc(\Theta_0,S,p) = \widehat p$. On a
morphism $(r,f) \colon (\Theta_0,S,p) \to (\Theta_0',S',p')$, a function
$r \colon \Theta_0 \to \Theta_0'$ and a measurable $f \colon S \to S'$ with
$\Prb(f) \circ p = p' \circ r$ (Definition~\ref{def:statmod}), put
$\Jc(r,f) = (\del r, \del f)$. Here $r$ is measurable because its domain is
discrete, so both components are Dirac kernels, hence deterministic
(Definition~\ref{def:dirac}); by Lemma~\ref{lem:square-pointwise} the pair is
a square from $\widehat p$ to $\widehat{p'}$ iff
$\widehat{p'}(r(\theta)) = f_{*}\,\widehat p(\theta)$ for every $\theta$,
i.e.\ iff $p' \circ r = \Prb(f) \circ p$, which is the hypothesis. Since
$\del$ is a functor and composition in $\Arrdet{\Stoch}$ is componentwise,
$\Jc$ preserves identities and composites.

\emph{Part~(1).} By the conventions above, $(\Theta_0,S,p) \mapsto
\widehat p$ is a bijection from the objects of $\StatMod$ onto the arrows of
$\Stoch$ whose domain carries its full power set, with inverse
$q \mapsto (A,S,\check q)$ for $q \colon A \cdot I \rightsquigarrow S$: the
domain of $\widehat p$ determines $\Theta_0$ as its underlying set, the
codomain determines $S$, and $p = \check{\widehat p}$. In particular $\Jc$ is
injective on objects and lands in $\Arrcop{\Stoch}$. For essential
surjectivity onto $\Arrcop{\Stoch}$, let $q \colon X \rightsquigarrow S$ with
$X \cong A \cdot I$ in $\Stoch$, say via mutually inverse
$v \colon A \cdot I \rightsquigarrow X$ and
$u \colon X \rightsquigarrow A \cdot I$. We show that both are deterministic,
arguing from the copower side, where singletons are measurable. For $a \in A$
put $X_a \coloneqq \{x \in X : u(x)(\{a\}) = 1\}$, measurable since
$x \mapsto u(x)(\{a\})$ is. Since $u \circ v = \id_{A \cdot I}$,
\[
1 \;=\; \delta_a(\{a\}) \;=\; (u \circ v)(a)(\{a\})
\;=\; \int_X u(x)(\{a\})\, v(a)(\mathrm{d}x),
\]
and the integrand is bounded by $1$, so $v(a)(X_a) = 1$; in particular
$X_a \neq \emptyset$. For $x \in X_a$ the probability measure $u(x)$ on
$2^{A}$ gives mass $1$ to $\{a\}$, so $u(x) = \delta_a$, and then, since
$v \circ u = \id_X$,
\[
\delta_x(F) \;=\; (v \circ u)(x)(F)
\;=\; \int_A v(b)(F)\, u(x)(\mathrm{d}b) \;=\; v(a)(F)
\qquad (F \in \Sigma_X).
\]
Thus $v(a) = \delta_x$ for every $x \in X_a$. Choosing $x_a \in X_a$ for each
$a$ (axiom of choice) gives a function $f \colon A \to X$, measurable because
its domain is discrete, with $v = \del f$; so $v$ is deterministic, and $u$
is deterministic by Lemma~\ref{lem:app-det}(4). Now
$q \circ v \colon A \cdot I \rightsquigarrow S$ has domain $A \cdot I$, so
$q \circ v = \Jc\bigl(A, S, (q \circ v)^{\vee}\bigr)$; and
$(v, \id_S) \colon q \circ v \to q$ and $(u, \id_S) \colon q \to q \circ v$
are squares, $q \circ v = \id_S \circ (q \circ v)$ and
$(q \circ v) \circ u = q = \id_S \circ q$, which are mutually inverse in
$\Arrdet{\Stoch}$, since $(u,\id_S) \circ (v,\id_S) = (u \circ v, \id_S) =
(\id, \id)$ and symmetrically. Hence
$q \cong \Jc\bigl(A,S,(q \circ v)^{\vee}\bigr)$.

\emph{Part~(2).} Fix $x = (\Theta_0,S,p)$ and $x' = (\Theta_0',S',p')$
with $x'$ satisfying~(a) and~(b), and let $(a,b)$ be a square from
$\widehat p$ to $\widehat{p'}$. Its left leg
$a \colon \Theta_0 \cdot I \rightsquigarrow \Theta_0' \cdot I$ is
deterministic, so by Lemma~\ref{lem:app-det}(1) each $a(\theta)$ is a
countably additive $\{0,1\}$-valued probability measure on $2^{\Theta_0'}$;
by~(a) it is $\delta_{r(\theta)}$ for a unique $r(\theta) \in \Theta_0'$, and
$r$ is measurable because $\Theta_0$ is discrete, so $a = \del r$. The right
leg is $b = \del f$ for a unique measurable $f \colon S \to S'$, by~(b) for
$S'$ and Lemma~\ref{lem:app-det}(3). By Lemma~\ref{lem:square-pointwise},
$(\del r, \del f)$ is a square iff $p' \circ r = \Prb(f) \circ p$, so
$(r,f)$ is a morphism of $\StatMod$ with $\Jc(r,f) = (a,b)$: the map
$\StatMod(x,x') \to \Arrdet{\Stoch}(\Jc x, \Jc x')$ is surjective. It is
injective by Lemma~\ref{lem:app-det}(2), since $2^{\Theta_0'}$ separates
points and $\Sigma_{S'}$ does by~(b).

\emph{Part~(3).} Write $\mathcal A$ for the full subcategory of
$\Arrcop{\Stoch}$ described in the statement. For $x = (\Theta_0,S,p)$ in
$\StatMod_{(a),(b)}$ the object $\Jc x = \widehat p$ has domain
$\Theta_0 \cdot I$ with $\Theta_0$ satisfying~(a) and codomain $S$
satisfying~(b), so $\Jc$ restricts to a functor
$\StatMod_{(a),(b)} \to \mathcal A$. It is fully faithful by part~(2),
because every object of $\StatMod_{(a),(b)}$ satisfies~(a) and~(b), and it
is essentially surjective by the argument of part~(1): for
$q \colon X \rightsquigarrow S$ in $\mathcal A$, with
$v \colon A \cdot I \rightsquigarrow X$ an isomorphism and $A$
satisfying~(a) as in the definition of $\mathcal A$, the object
$\bigl(A,S,(q \circ v)^{\vee}\bigr)$ lies in $\StatMod_{(a),(b)}$ and
$q \cong \Jc\bigl(A,S,(q \circ v)^{\vee}\bigr)$. A fully faithful and
essentially surjective functor is an
equivalence~\cite[IV.4, Thm.~1(iii)]{maclane_categories_1998}.

\emph{The finite case.} Every object of $\FinStoch$ is $A \cdot I$ for a
finite set $A$, so the object bijection of part~(1) restricts to a
bijection from the objects of $\StatMod_{\mathrm{fin}}$ onto the morphisms of
$\FinStoch$, i.e.\ onto the objects of $\Arrdet{\FinStoch}$ (which, $\FinStoch$
being full in $\Stoch$, is the full subcategory of $\Arrdet{\Stoch}$ on these
objects). Conditions~(a) and~(b) hold for every finite set and every finite
discrete space, so by part~(2) $\Jc$ is bijective on every hom-set. A
functor bijective on objects and on hom-sets is an isomorphism of categories,
so $\StatMod_{\mathrm{fin}} \cong \Arrdet{\FinStoch}$, which is more than the
stated equivalence.
\end{proof}

\bibliographystyle{plain}
\bibliography{calamus_bridge_draft_clean_v33k}

\end{document}
